\section{Historical v9 data used for the reported checkpoints}
\label{sec:data}

\subsection{Snapshot construction}

All checkpoint results reported in this manuscript use the historical
\texttt{combined\_dataset\_v9.npz} archive.  The parameterized replacement
dataset discussed below has not yet been used to train a reported model.
Each historical training example is a tuple
\begin{equation}
  (\eta_j,\xi_j,h_j)\longmapsto q_j^\star,
  \label{eq:data-map}
\end{equation}
on the $N=1024$ grid, where $q_j^\star$ is the stored reference target under
that source's generation protocol.  Initial conditions are sampled by wave
family and integrated with the double-precision reference method of
\cref{sec:solver}.  For rollout-derived sources we save multiple trajectory
times, recompute the order-six, pad-eight DNO label in manageable chunks, and
project the label to $|k|\leq128$.  Three static sources (IDs 2--4) instead
store their unfiltered order-six labels.  The mode loss uses only
$1\leq k\leq128$, where the conventions agree; the base loss also supervises
the stored high modes of the static sources.  Fields are stored as 32-bit
arrays, while trajectory generation and label evaluation use float64.

The current v9 dataset contains \num{7633989} snapshots and occupies
approximately \SI{88}{GB}.  Its composition is shown in
\cref{tab:data-composition}.  The families deliberately combine coherent
solutions, spectrally broad random fields, shallow and finite-depth cases, and
nonlinear modulation, broadening the training support beyond solitary-wave
states.

The archive uses fourteen source identifiers to record the generating shard.
In particular, suffixes such as \texttt{g0}, \texttt{g1}, and
\texttt{modal} identify independent generation shards or entry points.  They
do not denote different DNO truncation orders, target definitions, or physical
wave families.  We pool those shards in \cref{tab:data-composition}.

\begin{table}[t]
  \centering
  \caption{Construction-category composition of the historical
  \texttt{combined\_dataset\_v9.npz}, with independent generator shards pooled.}
  \label{tab:data-composition}
  \begin{tabular}{lr@{\qquad}lr}
    \toprule
    Construction category & Snapshots & Construction category & Snapshots\\
    \midrule
    Random sea       & 460,740   & Manufactured harmonic & 500,000\\
    Stokes           & 1,000,000 & Tanaka          & 2,200,000\\
    BF-inspired wavetrain & 2,033,249 & Sparse Fourier packets & 1,440,000\\
    \midrule
    \multicolumn{3}{r}{Total} & 7,633,989\\
    \bottomrule
  \end{tabular}
\end{table}

\subsection{Three roles of the historical construction categories}

The six historical construction categories in \cref{tab:data-composition}
are grouped below according to the states they contribute.  This grouping is
independent of the GPU or generator shard that produced a snapshot.

\paragraph{Coherent traveling and periodic waves (3,200,000 snapshots).}
This group contains both Tanaka and Stokes sources.  Single- and multi-crest
Tanaka conditions are constructed from numerical solitary-wave profiles, with
crest amplitude sampled relative to depth and positions randomized subject to
a minimum separation.  The main shallow set reaches $h=0.01$ and uses an
occupancy-aware upper bound to avoid cutting a broad tail at the periodic
boundary.  The steep addition contributes 200,000 snapshots in
$h\in[0.20,0.35]$ and $a/h\in[0.25,0.45]$.  Deep- and finite-depth Stokes waves
supply periodic states with controlled harmonic content.  Together these
sources resolve coherent nonlinear structure across depth without treating
the deep/finite data partitions as separate physical families.

\paragraph{JCP09-inspired modulated wavetrains (2,033,249 snapshots).}
The historical constructor starts from a fifth-order Stokes carrier and two
sidebands inspired by \cite{xu2009}, but it is not the literal JCP09 initial
condition.  It uses one carrier coefficient for both sidebands, adds empirical
cross terms, and applies the deep-water sideband formula at sampled finite
depths.  Only \(46.7\%\) of the archived parameter draws lie inside the
leading deep-water modulational-instability band.  The three archived shards
are independent samples from this same construction; the suffix
\texttt{modal} does not define a separate family.

\paragraph{Spectral-coverage and perturbation states (2,400,740 snapshots).}
This group combines the manufactured-harmonic, random-sea, and three sparse
Fourier-packet sources.  Manufactured harmonics anchor the small-amplitude
limit, while random seas broaden phase, bandwidth, and local-slope coverage.
The sparse packet sources address the
neighborhood of solitary waves rather than the exact solitary-wave manifold:
they combine one to five Fourier modes while controlling $kh$, $a/h$, and
$ka$.

\subsection{Forward paper-dataset design}

The definitions in this subsection belong to a replacement dataset that has
not yet been used for training.  In particular, they do not describe a
retrospective cleaning of v9 and cannot explain the errors of the checkpoints
reported in this manuscript.  The replacement has four physical families:
fifth-order Stokes states, periodic tangent-Hermite Tanaka profiles, the
equation-(33) sideband construction paired with the project's analytic
fifth-order deep-water carrier, and resolved-band JONSWAP/TMA random seas.
The four physical constructors and their parameter samplers are implemented.
The current family revisions are 2 for Stokes, 3 for Tanaka, and 4 for both
Benjamin--Feir and JONSWAP/TMA.  A revision number is therefore local
to a physical family, rather than a single dataset-wide version number.
Unless an endpoint convention is written explicitly, closed brackets in a
displayed parameter box denote the closure of the continuous support.  The
ordinary floating-point uniform draw uses its lower endpoint and excludes its
upper endpoint.  An endpoint convention written explicitly in an individual
sampling formula---including the half-open phase and center intervals and the
mixed open--closed Benjamin--Feir steepness draw---overrides that default.
These endpoint choices do not change a continuous uniform law, but they are
specified for deterministic regeneration.
For each trajectory family, one common executor runs a
residual-controlled GL2 trajectory with time step \(0.01\), saving a
candidate state every \(0.08\).  JONSWAP/TMA first uses the explicitly
time-dependent adjustment defined below and then starts its autonomous
production trajectory.  A case is retained only when its declared numerical
path returns a complete admissible trajectory on the requested interval.  The
\(0.005\) and \(0.0025\) steps are used only by a separate method-validation
utility; they are not additional production rollouts or retry steps.

Common modules assign a split and parameter category before numerical work and
write one atomic completed-batch artifact containing the family and split, the
ordered parameter-group identifiers, and the accepted row arrays.  A batch slot
with no owned row block denotes a rejected attempt; sampled specifications and
phases, failure reasons, and diagnostic metrics are not persisted.  A resumable
driver continues until the accepted-case quota in every parameter category is
met, replacing a rejected attempt only in its original category.  The fixed
split, family, and attempted-simulation index seed deterministic regeneration.
The training loader uses the preassigned splits and globally shuffles all
accepted training rows.  Unexpected construction or target-evaluation errors
propagate rather than being silently converted into ordinary rejected samples.
The family laws below therefore describe attempted specifications, not the
unconditional law of loader-visible cases.  If \(X\) is a complete attempted
specification in preassigned category \(i\), \(A_f\) is the declared complete
case-acceptance event for family \(f\), and \(w_i\) is the accepted-quota
weight, then the released law is
\begin{equation}
  p_{\mathrm{release}}(i,x)
  =w_i\,p_{\mathrm{proposal}}(x\mid i,A_f=1).
  \label{eq:paper-accepted-law}
\end{equation}
For Stokes, \(A_f\) comprises declared static support, a finite state, positive
water column, and a finite target.  For Tanaka, it comprises declared
construction support followed by complete autonomous-trajectory acceptance;
for Benjamin--Feir, initial-state validity followed by complete autonomous-
trajectory acceptance.  For JONSWAP/TMA, it comprises initial construction and
graph validity followed by nonlinear adjustment and complete autonomous-
trajectory acceptance.  Thus the quotas preserve the declared category
marginals, not the unconditioned proposal law within a category.  This is a
statement about how the release is constructed; it does not by itself assert a
measured within-category shift for every family.  No retrospective parameter
cutoff is applied.  The loader-facing manifest and row map record the selected
shards, split totals, and case-to-row ownership; they do not persist numerical-
configuration or source hashes.  The release described here selects Stokes
revision 2, Tanaka revision 3, Benjamin--Feir revision 4, and JONSWAP/TMA
revision 4.

An earlier generator-revision-1 complete transaction smoke test used a reduced
wiring contract:
64 spatial grid points, Craig--Sulem order zero, padding factor one, largest
delivered wavenumber 16, and the former paired \(0.04/0.02\) validation path.
One Tanaka case, one Benjamin--Feir case, and one random-sea case were each
accepted to a quota of one case; each committed three rows, and all nine rows
were finite when read by the training loader.  This reduced calculation tests
the transaction and loader, not the current family-revision executors, the
paper target, the long-horizon numerical contract, or a population acceptance
rate.  The current numerical paths have focused software-test coverage.
A former revision-2 104-category trajectory pilot predates the Tanaka
revision-3 support and numerical contract and remains historical software-path
evidence.  The present family specifications contain 108 categories: four
Stokes, eleven Tanaka, 66 Benjamin--Feir, and 27 JONSWAP/TMA.  The exact
revision-3 Tanaka pilot passed all eleven Tanaka categories.  The revision-4
Benjamin--Feir exact-source gate passed all 66 categories, and its production
splits are complete at 16384 training cases and 1024 cases in each of
validation and test.  The revision-4 JONSWAP/TMA source was tested more
broadly: a predeclared gate retained 540 trajectories, 20 in each of its 27
categories, from 548 attempts.  Its observed rejection rate was therefore
\(8/548=1.46\%\), below the predeclared two-percent bound.  This was the
completed constructor-adoption criterion, not a bulk release threshold.  The
subsequent revision-4 bulk release is complete: it retained \num{18432} cases
from \num{18726} attempts, recorded \num{294} zero-row rejections, and stored
\num{294912} rows.

The revision-2 Stokes splits are complete at the same accepted-case counts as
Benjamin--Feir.  A historical Tanaka run produced 2048 training, 1024
validation, and 1024 test cases, but those trajectories are excluded from the
paper dataset.  The old solitary-wave solve could not realize crest amplitudes
below \(9.9478\times10^{-4}\); 64 of the 4096 historical cases contain one
requested component below this floor.  The smallest requested component was
inflated by a factor \(41.27\) after periodic construction and projection.
The corrected solve uses \(q_c^{\max}=1-10^{-12}\), 48 bisections, and checks
that every realized crest height matches its request.  The corrected
revision-3 release is a fresh deterministic population with 16384
training cases and 1024 cases in each of validation and test.  It does not
reuse the historical random specifications, and no old and corrected Tanaka
rows are mixed.  All \num{18432} corrected attempts were accepted and stored
\num{3686400} rows.

\paragraph{One frozen learning target.}
Let the spatial period be $L=2\pi$, let $N=1024$, and let
$\Delta x=L/N$ and $x_\ell=\ell\Delta x$, $\ell=0,\ldots,N-1$, be the grid
spacing and nodes.  For a periodic grid function $u$, let $P_K u$ retain its
Fourier modes with $|k|\leq K$, where $K=128$, and let $P_K^\circ u$ also set
the zero Fourier mode to zero.  We
write $\G_M^{N,p}(\eta;h)\xi$ for the order-$M$ Craig--Sulem approximation on
$N$ points with product-padding factor $p$, surface elevation $\eta$, surface
potential $\xi$, and depth $h>0$.  The target is fixed to
\begin{equation}
  q_{\mathrm{ref}}(\eta,\xi;h)
  =
  P_K^\circ\!\left[
    \G_{6}^{1024,8}(P_K\eta;h)(P_K\xi)
  \right].
  \label{eq:forward-dataset-target}
\end{equation}
Thus both inputs are restricted before evaluation, and the output is
restricted and made mean zero afterward.  The target is evaluated in float64;
the three fields $\eta$, $\xi$, and $q_{\mathrm{ref}}$ are stored in float32,
while depth and retained time are stored in float64.  The sampled parameters
from which a row was constructed are not stored.
Thus the mean-zero statement is exact for the float64 target and is preserved
in the stored array only to float32 rounding error.
Static Stokes states use exactly \cref{eq:forward-dataset-target}, rather than
the unfiltered labels in some historical v9 sources.  The paper generator fixes
$L,N,K,M,p$ and the evaluation precision in its numerical settings; the completed
batches do not record target-configuration or source hashes.
The three trajectory families have internal evolution contracts specified
below, but all deliver the same target in
\cref{eq:forward-dataset-target}.

JONSWAP/TMA is evolved on a larger grid than the stored target.  To specify
the transfer without an FFT-normalization ambiguity, for a field sampled on
\(N'\) points define
\begin{equation}
  \widehat f_k^{(N')}
  =\frac1{N'}\sum_{j=0}^{N'-1}
    f\!\left(\frac{2\pi j}{N'}\right)e^{-2\pi i jk/N'}.
  \label{eq:paper-normalized-fourier-coefficient}
\end{equation}
At each retained JONSWAP/TMA time, the coefficients of \(\eta\) and \(\xi\)
with \(|k|\le128\) are copied from the 2048-point internal grid and
reconstructed on the 1024-point grid.  Equation
\eqref{eq:forward-dataset-target} is then evaluated from those reconstructed
fields.  The internal value of \(G(\eta;h)\xi\) is never transferred between
grids; recomputing it is necessary because the stored label is the declared
1024-point order-six operator, not the order-four operator used for evolution.

\paragraph{Static fifth-order Stokes states.}
Let $n\geq1$ be an integer carrier mode, $k=2\pi n/L$ its wavenumber,
$a>0$ the amplitude parameter in Fenton's fifth-order expansion
\cite{fenton1985}, $h>0$ the depth, and
$\phi\in[0,2\pi)$ a spatial phase.  The finite- and deep-water coefficient
tables are two parameter branches of the same fifth-order family.  In either
branch the constructor evaluates the state directly at
\begin{equation}
  \theta=kx+\phi,
  \qquad \phi\sim\operatorname{Unif}[0,2\pi).
  \label{eq:paper-stokes-phase}
\end{equation}
No sampled time is used as a substitute for $\phi$: changing $a$ changes the
nonlinear temporal frequency, whereas \cref{eq:paper-stokes-phase} preserves
the requested spatial translation.  The spatial mean of $\xi$ is removed to
fix the irrelevant constant-potential gauge.

The common amplitude interval is
\begin{equation}
  0.000766\leq a\leq0.011494.
  \label{eq:paper-stokes-amplitude}
\end{equation}
In the finite-depth branch,
\begin{equation}
  n\in\{14,\ldots,26\},\qquad
  0.02\leq h\leq1.5,\qquad
  0.5\leq kh\leq5,
  \label{eq:paper-stokes-finite-support}
\end{equation}
whereas the deep-water branch uses
\begin{equation}
  n\in\{1,\ldots,20\},\qquad
  4\leq h\leq50,\qquad
  kh\geq5.
  \label{eq:paper-stokes-deep-support}
\end{equation}
The steepness $ka$ is assigned to one of the two categories
\begin{equation}
  0.005\leq ka<0.03,
  \qquad
  0.03\leq ka\leq0.15.
  \label{eq:paper-stokes-cells}
\end{equation}
The finite/deep branches crossed with these two intervals form four parameter
categories.  Within an assigned category, the sampler first enumerates the carrier
modes for which the branch, depth, amplitude, and steepness intervals have
nonempty intersection, and chooses one of those modes uniformly.  It then
draws \(h\) log-uniformly on the mode-conditioned depth interval, draws
\(\phi\) uniformly on \([0,2\pi)\), and draws \(a\) uniformly from the
intersection of \cref{eq:paper-stokes-amplitude} with the assigned steepness
interval.

The finite-depth branch has an additional support condition.  Write its
elevation as
\begin{equation}
  \eta_5(\theta)=\sum_{m=1}^{5}E_m\cos(m\theta),
  \label{eq:paper-stokes-harmonics}
\end{equation}
where $E_m$ is the coefficient of the $m$th harmonic.  Let
$H=\max_\theta\eta_5(\theta)-\min_\theta\eta_5(\theta)$ be the exact
crest-to-trough height, and let $H_{4096}$ be the corresponding maximum minus
minimum over 4096 equally spaced values of $\theta$.  Define
\begin{equation}
  H_+
  =
  H_{4096}
  +\frac{1}{4}\left(\frac{2\pi}{4096}\right)^2
    \sum_{m=1}^{5}m^2|E_m|,
  \qquad
  \mathrm{Ur}_+
  =
  \frac{H_+\lambda^2}{h^3},
  \qquad
  \lambda=\frac{2\pi}{k}.
  \label{eq:paper-stokes-ursell}
\end{equation}
Since $|\eta_5''|\leq\sum_{m=1}^5m^2|E_m|$, Taylor's theorem at the true
maximum and minimum gives $H\leq H_+$.  The finite-depth sampler therefore
requires the conservative condition $\mathrm{Ur}_+\leq26$.  This is a
restriction on the parameter support of the fifth-order family, not a filter
trained on integration failures.  Fenton identifies the corresponding
shallow-water expansion parameter for fifth-order Stokes theory
\cite{fenton1985}, and Zhao, Wang, and Liu use $\mathrm{Ur}=26$ to separate
the fifth-order Stokes and shallower cnoidal regimes \cite{zhao2024}.
The coefficient implementation follows the fifth-order formula checked
against Zhao and Liu \cite{zhao2022}.  The deep-water branch instead requires
$kh\geq5$.

\paragraph{Seeded sampling.}
Let $s$ be the root seed fixed by the data split, let $f$ identify the physical
family, and let $i$ be the attempted-simulation index.  A simulation uses NumPy
PCG64 initialized from $(s,f,i)$.  The four
categories are the finite/deep branches crossed with the two intervals in
\cref{eq:paper-stokes-cells}.  Requested accepted-case counts are divided
among these categories by quotient and remainder before any state is constructed.
This is a seeded stratified pseudorandom construction, not a low-discrepancy
sequence.  A successful proposal holds its realized mode, amplitude, depth,
phase, branch, and category only long enough to construct its rows; these
values can instead be regenerated from the split, family, and attempt seed.

If a finite-depth amplitude draw violates
$\mathrm{Ur}_+\leq26$, only the amplitude is redrawn, uniformly from the same
category intersection.  Its mode, depth, phase, and category remain fixed.  The case
may make at most 1000 additional amplitude draws.  Intermediate rejected
amplitudes and their values of $\mathrm{Ur}_+$ are not persisted.  If the
redraw allowance is exhausted, the sampler returns no state and the batch slot
owns no rows; generation then fills the remaining
accepted-case quota from a later attempt in the same category.

\paragraph{Static archive and replay checks.}
These checks concern only the implemented Stokes component of the replacement
dataset.  A 16-case finite-depth archive contained eight cases from each category
in \cref{eq:paper-stokes-cells}, required three same-category redraws, and had
$\max\mathrm{Ur}_+=18.0578589$.  A separate 16-case deep archive was finite
and had $\min kh=10.0678$.  Its category counts were nine and seven because some
carrier modes admitted only one steepness category; the subset admitting both remained
balanced to within one.  Re-evaluating
\cref{eq:forward-dataset-target} from the stored float64 parameters and then
converting the fields to float32 reproduced both archives bit for bit.  A
synthetic interruption after two batches likewise resumed to the same 30
ordered archive entries, byte for byte, as an uninterrupted 32-case run.  In
a 20,000-case finite-depth sampling stress with root seed 20260725 and stream
zero, the two category counts were 10,000 and 10,000, the largest accepted value
was $\mathrm{Ur}_+=25.9724665$, the largest redraw count was 56, and no case
exhausted the cap.  A separate forced-exhaustion quota test committed the
attempt as a zero-row out-of-support rejection and scheduled a
replacement in the same category.

The common static transaction was also exercised at the exact target
resolution $(N,M,p,K)=(1024,6,8,128)$ in float64.  One validation attempt was
assigned to each finite/deep and low/moderate category, and its complete
specification was written before construction.  All four attempts had finite
input and target fields, positive water column, no support violation, and no
failed required quality check.  The transaction therefore committed one
$t=0$ row per case and constructed the split-aware dataset view.  This
accepted-quota calculation accepted all four attempts and checks the exact
implementation path; it does not estimate a population acceptance
probability.

\paragraph{Periodic tangent-Hermite Tanaka states.}
Let $\alpha=a/h>0$ be the ratio of crest amplitude $a$ to depth $h$.
Following Tanaka's solitary-wave formulation \cite{tanaka1986}, a unit-depth
solve at amplitude ratio $\alpha$ returns a dimensionless surface profile
$\bar\eta_\alpha(X)$, a tabulated surface angle $\theta_\alpha(X)$, and a
Froude number $F_\alpha$.  Xu and Guyenne also use a solitary wave computed
by Tanaka's method \cite{xu2009}.  The tangent-Hermite interpolation,
periodization, multi-crest superposition, and parameter sampling below are
choices of the present dataset; they are not parts of either cited
solitary-wave construction.  Here $X$ is the dimensionless horizontal
coordinate, and
\begin{equation}
  \frac{\dd\bar\eta_\alpha}{\dd X}(X)
  =\tan\theta_\alpha(X).
  \label{eq:paper-tanaka-slope}
\end{equation}
Before interpolation, the constructor verifies that both endpoint elevations
and the endpoint and crest tangents differ from zero by at most $10^{-12}$,
then sets those quantities exactly to zero.  The corrected tabulated profile
is extended by zero outside its finite tabulation interval.  Between adjacent
tabulation points, the constructor uses the unique cubic polynomial matching
both endpoint values of $\bar\eta_\alpha$ and both endpoint values of
$\tan\theta_\alpha$.  Thus ``tangent-Hermite'' specifies the interpolation,
not an additional wave model.  The physical derivative is correctly scaled
because
\begin{equation}
  \partial_x\!\left[
    h\bar\eta_\alpha\!\left(\frac{x-x_0+rL}{h}\right)
  \right]
  =
  \frac{\dd\bar\eta_\alpha}{\dd X}
  \!\left(\frac{x-x_0+rL}{h}\right).
  \label{eq:paper-tanaka-scaled-slope}
\end{equation}

For a crest centered at $x_0\in[0,L)$, define the periodic elevation
\begin{equation}
  \eta_{\alpha,h,x_0}(x)
  =
  h\sum_{r\in\mathbb Z}
  \bar\eta_\alpha\!\left(\frac{x-x_0+rL}{h}\right).
  \label{eq:paper-tanaka-periodization}
\end{equation}
Only finitely many terms in \cref{eq:paper-tanaka-periodization} are nonzero
at a given $x$, because the extended profile has finite support.  The
implementation sums every periodic image whose support intersects $[0,L)$;
it does not retain only the nearest three images.  The sum is evaluated on an
$8N$-point grid and Fourier-restricted to the $N$-point grid.

Let $g>0$ be gravitational acceleration, let
$c_\alpha=F_\alpha\sqrt{gh}$ be the unsigned wave speed, and let
$s\in\{-1,1\}$ denote the direction of travel.  Write
$\eta_x=\partial_x\eta_{\alpha,h,x_0}$.  The traveling-wave relation first
defines
\begin{equation}
  r_{\alpha,h,x_0,s}(x)
  =
  s\left[
    c_\alpha-
    \sqrt{(1+\eta_x(x)^2)
      (c_\alpha^2-2g\eta_{\alpha,h,x_0}(x))}
  \right].
  \label{eq:paper-tanaka-potential-derivative}
\end{equation}
Define the expression under the square root by
\begin{equation}
  \mathcal R_{\alpha,h,x_0}(x)
  =
  (1+\eta_x(x)^2)
  (c_\alpha^2-2g\eta_{\alpha,h,x_0}(x)).
  \label{eq:paper-tanaka-radicand}
\end{equation}
The constructor requires the nodewise condition
$\min_\ell\mathcal R_{\alpha,h,x_0}(x_\ell)\geq0$ and treats an attempted component
that violates this condition as inadmissible; it does not replace a negative
value by zero.  For the fixed upper-steepness, seam-centered test case, the
initial value is
$\min_\ell\mathcal R_{\alpha,h,x_0}(x_\ell)/c_\alpha^2=0.372151$.  Hence the clipping
in the old writer did not change that particular initial condition.  This
single positive margin does not establish the condition for the full
population.

Let $\widehat r_n$ be the Fourier coefficient of
$r_{\alpha,h,x_0,s}$ at wavenumber $k_n=2\pi n/L$.  The periodic,
zero-mean surface potential is defined by
\begin{equation}
  \widehat\xi_n
  =
  \begin{cases}
    \widehat r_n/(ik_n),&n\ne0,\\
    0,&n=0.
  \end{cases}
  \label{eq:paper-tanaka-potential}
\end{equation}
Consequently $\partial_x\xi=r_{\alpha,h,x_0,s}-\widehat r_0$; subtracting
the constant $\widehat r_0$ is necessary for a periodic antiderivative.

A case with $m\geq1$ crests has parameters
$(h,m,\alpha_j,x_j,s_j)_{j=1}^m$ and initial state
\begin{equation}
  (\eta_0,\xi_0)
  =
  P_K\sum_{j=1}^{m}
  \left(
    \eta_{\alpha_j,h,x_j},
    \xi_{\alpha_j,h,x_j,s_j}
  \right).
  \label{eq:paper-tanaka-state}
\end{equation}
For $m>1$, \cref{eq:paper-tanaka-state} is a superposition of separately
constructed solitary-wave profiles, not an exact steady multi-solitary-wave
solution.
The sharp $P_K=P_{128}$ map in \cref{eq:paper-tanaka-state} is applied once
before time integration.  Revision 3 then embeds this state in the larger
$|k|\leq256$ evolution space by setting modes $129\leq|k|\leq256$ initially
to zero.  Nonlinear evolution may populate those workspace modes, but the
initial-value problem is exactly the sharp-$P_{128}$ problem used in the
resolution calculation below.
For two centers $x_i,x_j\in[0,L)$, define their periodic distance by
\begin{equation}
  d_L(x_i,x_j)
  =
  \min_{r\in\mathbb Z}|x_i-x_j+rL|.
  \label{eq:paper-periodic-distance}
\end{equation}
The centers must satisfy $d_L(x_i,x_j)\geq3h$ whenever $i\ne j$.

The amplitudes are sampled before the depth.  In a main group, define the
total dimensionless amplitude
$\alpha_\Sigma=\sum_{j=1}^m\alpha_j$ and draw
$\alpha_\Sigma\sim\operatorname{Unif}[0.10,0.35]$.  Draw nonnegative weights
$(\lambda_1,\ldots,\lambda_m)$ uniformly on the simplex
$\sum_j\lambda_j=1$, equivalently from
$\operatorname{Dirichlet}(1,\ldots,1)$, and set
$\alpha_j=\alpha_\Sigma\lambda_j$.  For $m=1$, this means
$\lambda_1=1$.  In a steep group, draw the single amplitude ratio
$\alpha_1\sim\operatorname{Unif}[0.25,0.45]$.  In both groups, define
\begin{equation}
  \alpha_{\max}=\max_{1\leq j\leq m}\alpha_j.
  \label{eq:paper-tanaka-maximum-amplitude}
\end{equation}

The depth interval is then restricted before a profile is constructed.  The
small-amplitude KdV profile has the form
\begin{equation}
  \frac{\eta(x)}{h}
  \approx
  \alpha_{\max}\operatorname{sech}^2\!\bigl(\kappa(x-x_0)\bigr),
  \qquad
  \kappa=\frac{\sqrt{3\alpha_{\max}}}{2h},
  \label{eq:paper-tanaka-kdv-width}
\end{equation}
where $\operatorname{sech}z=1/\cosh z$ and $\kappa$ is an inverse length.
We use $\kappa$ only as a small-amplitude KdV inverse-width proxy for the
narrowest component.  It is not the exact width of a Tanaka profile.  With
the delivered cutoff $K=128$, define the number of delivered wavenumbers per
proxy inverse width by
\begin{equation}
  R=\frac{128}{\kappa}
   =\frac{256h}{\sqrt{3\alpha_{\max}}}.
  \label{eq:paper-tanaka-resolution-ratio}
\end{equation}
Tanaka revision 3 requires $R\geq10$.  This threshold is a numerical design
choice established by the resolution tests below; it is not a condition
stated in \cite{tanaka1986} or \cite{xu2009}.

The proxy first arose from a retrospective resolution audit, rather than from
a fitted rejection statistic.  We reconstructed 256 tangent-Hermite proposals
before their sharp $P_{128}$ projection.  All eleven cases that exceeded the
historical $K_{\mathrm{evol}}=224$ Hamiltonian-drift diagnostic had $R<8$.
The 210 cases with $R\geq8$ contained none of those diagnostic failures, and
their largest raw-to-$P_{128}$ relative root-mean-square errors were
$6.80\times10^{-5}$ for $\eta$ and $4.47\times10^{-6}$ for $\xi$.  By
comparison, the unrelated fixed restriction $h\geq0.025$ admitted a larger
maximum $\eta$ projection error, $4.33\times10^{-4}$.  These observations
identify under-resolved input width as the mechanism; the full-horizon panel
below, not the historical Hamiltonian threshold, determines the revision-3
support.

For a base lower depth $b>0$, define
\begin{equation}
  h_-(\alpha;b)
  =\max\!\left\{b,\frac{10\sqrt{3\alpha}}{256}\right\}.
  \label{eq:paper-tanaka-depth-lower-bound}
\end{equation}
After the amplitudes have been drawn, sample $U_h\sim\operatorname{Unif}[0,1]$
and set
\begin{equation}
  h=h_-^{\,1-U_h}h_+^{\,U_h},
  \label{eq:paper-tanaka-geometric-depth}
\end{equation}
where $h_-=h_-(\alpha_{\max};0.01)$ and $h_+=0.30$ in a main group, while
$h_-=h_-(\alpha_1;0.20)$ and $h_+=0.35$ in a steep group.  Thus $h$ is
log-uniform on the conditional interval $[h_-,h_+]$.  The resolution term in
\cref{eq:paper-tanaka-depth-lower-bound} is at most $0.0454$ on the steep
amplitude interval, so the steep lower bound remains $0.20$ and the
resolution restriction is inactive there.  The resulting supports are
therefore
\begin{align}
  &m\in\{1,2,3\},\quad
    \alpha_\Sigma\in[0.10,0.35],\quad
    h\in[h_-(\alpha_{\max};0.01),0.30]
    &&\text{(main)},
  \label{eq:paper-tanaka-main-support}\\
  &m=1,\quad
    \alpha_1\in[0.25,0.45],\quad
    h\in[0.20,0.35]
    &&\text{(steep)}.
  \label{eq:paper-tanaka-steep-support}
\end{align}
Equal increments of $U_h$ give equal multiplicative changes of depth.  This is
a coverage rule for the numerical experiment, not a probability model for
physical water depths.  In a deterministic million-draw audit for each
$m=1,2,3$, the $R\geq10$ condition removed $28.55\%$ of the former
independent-depth law after weighting the nine main categories by their
$2:3:4$ category counts.  The median depths for $m=1,2,3$ were respectively
$0.0973$, $0.0902$, and $0.0857$, and their 95th percentiles were $0.2680$,
$0.2659$, and $0.2647$.

For each case let
\begin{equation}
  q=\#\{j:s_j=+1\}
  \label{eq:paper-tanaka-right-count}
\end{equation}
be the number of right-moving crests.  The three choices of $m$ in
\cref{eq:paper-tanaka-main-support}, crossed with
$q=0,\ldots,m$, give nine parameter categories.  The steep support in
\cref{eq:paper-tanaka-steep-support}, crossed with $q=0,1$, gives two more.
Thus the Tanaka family has eleven parameter categories.

The requested Tanaka accepted-case count is divided equally across these
eleven categories, up to the unavoidable difference of one case when the
total is not divisible by eleven.  Conditional on $q$, the sampler chooses
uniformly which $q$ of the $m$ exchangeable crest labels have sign $+1$; all
remaining signs are $-1$.  Hence left and right directions remain balanced
marginally in the ideal equal-category population, while the direction
compositions $q=0,\ldots,m$ receive equal coverage instead of their former
binomial frequencies under independent fair signs.  For a finite quota, the
quotient--remainder assignment changes category counts by at most one and may
therefore introduce a one-case direction imbalance.  In particular, the
fixed category order gives a quota of 1024 one extra left-moving one-crest
case.  This law also gives more aggregate Tanaka weight to
$m=2$ and $m=3$, because they have three and four distinct direction
compositions, respectively.  The corrected full-size Tanaka production is
complete.  Its historical 4096-case predecessor and its eleven-case
revision-3 pilot are excluded from the replacement-dataset view.

Tanaka revision 3 retains the eleven direction-aware categories introduced in
revision 2, but changes the amplitude--depth draw order, the admissible depth
support, and the internal numerical contract.  Earlier Tanaka revision-1 and
revision-2 smokes, dry runs, and pilot artifacts remain historical
constructor or transaction evidence; they cannot be resumed into a
revision-3 run.  Benjamin--Feir and JONSWAP/TMA use revision 4; static Stokes
remains at revision 2.

To sample centers without a translation bias, let
$(w_1,\ldots,w_m)$ have the same uniform-simplex law and set the cyclic gaps
to
\begin{equation}
  g_j=3h+(L-3mh)w_j,\qquad j=1,\ldots,m.
  \label{eq:paper-tanaka-gap-law}
\end{equation}
A uniform global rotation places the resulting point set on the periodic
domain.  Directions are then assigned according to the category as described
above.  In a steep group, the center uses this same uniform rotation and the
direction is fixed by $q\in\{0,1\}$.  The separation is always feasible on
the declared support because $3mh\leq2.7<2\pi=L$.  Thus separation and
translation symmetry are enforced before integration.

The population sampler and ordered batch adapter are implemented.  The
sampled parameters are held only during construction and can be regenerated
from the split, family, and attempt seed.  The potential routine
transfers the computed radicand to the host once before taking the square root
and introduces no clipping tolerance.  An everywhere-finite, positive-speed
component with a negative radicand is a declared support rejection: it
contributes no rows and is replaced in the same category.  A nonfinite radicand,
a nonpositive or nonfinite computed speed, or an unexpected exception remains
an implementation error and stops the run.
In the historical reduced paired calculation described above, the Tanaka
attempt was accepted, committed three finite rows, and was finite when read by
the training loader.  This is transaction evidence at the reduced contract,
not evidence about the current Tanaka trajectory executor or the Tanaka
population.

Tanaka revision 3 evolves an internal band
$|k|\leq K_{\mathrm{evol}}$ with $K_{\mathrm{evol}}=256$.  Let
$\widehat\eta_{0,k}$ and $\widehat\xi_{0,k}$ be the Fourier coefficients of
the sharp-$P_{128}$ state in \cref{eq:paper-tanaka-state}.  The internal state
is initialized by the zero-filled lift
\begin{equation}
  \bigl(\widehat\eta_k(0),\widehat\xi_k(0)\bigr)
  =
  \begin{cases}
    \bigl(\widehat\eta_{0,k},\widehat\xi_{0,k}\bigr),&|k|\leq128,\\
    (0,0),&128<|k|\leq256.
  \end{cases}
  \label{eq:paper-tanaka-zero-filled-lift}
\end{equation}
Thus no discarded raw-profile mode is supplied to the wider evolution at
$t=0$; modes above 128 can arise only from the subsequent dynamics.  The
$K_{\mathrm{evol}}=224$ comparison arm uses the same construction with upper
limit 224.  After every completed GL2 step, the Fourier coefficients
$\widehat\eta_k$ and $\widehat\xi_k$ are multiplied by
\begin{equation}
  \sigma_k
  =\exp\!\left[-36\left(\frac{|k|}{256}\right)^{36}\right].
  \label{eq:paper-tanaka-hou-li-filter}
\end{equation}
This is the power-36 Hou--Li filter \cite{hou2007} used on both canonical
variables, with the parameters of \cite[Eq.~(29)]{xu2009}.  The
$K_{\mathrm{evol}}=224$ comparison arm replaces 256 in
\cref{eq:paper-tanaka-hou-li-filter} by 224.  At each saved time, the delivered
state is
$(P_{128}\eta,P_{128}\xi)$ and the delivered label is
$q_{\mathrm{ref}}$ from \cref{eq:forward-dataset-target}.  Modes
$129\leq|k|\leq256$ therefore stabilize the internal evolution but are not
stored as inputs or labels for learning.
This sharp \(P_{128}\) initial projection, zero-filled lift to
\(K_{\mathrm{evol}}=256\), post-step Hou--Li multiplier, and final
\(P_{128}\) delivery is the exact revision-3 contract used by the pilot.  The
excluded historical July training shards used the same numerical contract but
are not part of the fresh paper dataset.

The resolution choice was tested on nine one-crest cases.  Their amplitude
ratios were $\alpha\in\{0.225,0.30,0.35\}$, and their resolution ratios were
$R\in\{6,8,10\}$.  Each case was evolved from the same bit-identical
$P_{128}$ initial state with
$K_{\mathrm{evol}}=224$ and 256 through $T=200$.  Define the dimensionless
delivered fields
\begin{equation}
  \widetilde\eta_J(t)=\frac{P_{128}\eta_J(t)}{h},\qquad
  \widetilde\xi_J(t)=\frac{P_{128}\xi_J(t)}{h\sqrt{gh}},\qquad
  \widetilde q_J(t)
  =\frac{q_{\mathrm{ref}}(\eta_J(t),\xi_J(t);h)}{\sqrt{gh}},
  \label{eq:paper-tanaka-dimensionless-comparison-fields}
\end{equation}
where $J\in\{224,256\}$ is the internal cutoff and
$(\eta_J,\xi_J)$ is that arm's internal state.  For
$u\in\{\widetilde\eta,\widetilde\xi,\widetilde q\}$, define the reported
cross-resolution error by
\begin{equation}
  \epsilon_u
  =\max_t
  \frac{\lVert u_{224}(t)-u_{256}(t)\rVert_{2,N}}
       {\max_s\lVert u_{256}(s)\rVert_{2,N}+10^{-12}},
  \qquad
  \lVert v\rVert_{2,N}^2
  =\frac{1}{N}\sum_{\ell=0}^{N-1}|v(x_\ell)|^2.
  \label{eq:paper-tanaka-cross-resolution-error}
\end{equation}
The three fields are nondimensionalized before the additive $10^{-12}$ floor
is applied.  Write $\epsilon_\eta$, $\epsilon_\xi$, and $\epsilon_q$ for
\cref{eq:paper-tanaka-cross-resolution-error} evaluated with
$u=\widetilde\eta$, $\widetilde\xi$, and $\widetilde q$, respectively.  The
audit also checked the analogous full-history RMS error.  To measure whether
hidden state modes alter the delivered operator, define
\begin{equation}
  d_J(t)
  =P_{128}^\circ\!\left[
    \G_6^{1024,8}(\eta_J(t);h)\xi_J(t)
  \right]
  -q_{\mathrm{ref}}(\eta_J(t),\xi_J(t);h),
  \label{eq:paper-tanaka-hidden-band-defect}
\end{equation}
where $(\eta_J,\xi_J)$ is the internal state of arm $J$.  Its reported closure
error is
\begin{equation}
  C_J
  =\max_t
  \frac{\lVert d_J(t)\rVert_{2,N}}
       {\max_s\lVert q_{\mathrm{ref}}(\eta_J(s),\xi_J(s);h)\rVert_{2,N}
        +10^{-12}}.
  \label{eq:paper-tanaka-hidden-band-closure}
\end{equation}

All three $R=6$ controls failed the predeclared $10^{-3}$ comparison; their
worst operator error was $\epsilon_q=9.36\times10^{-3}$ and
their worst closure error was $2.91\times10^{-3}$.  All three $R=8$ cases
passed, but their worst errors were
$(\epsilon_\eta,\epsilon_\xi,\epsilon_q)
=(1.87\times10^{-4},1.29\times10^{-5},7.34\times10^{-4})$ and
$\max_J C_J=6.31\times10^{-4}$.  At $R=10$, these values fell to
$(2.94\times10^{-5},1.62\times10^{-6},1.35\times10^{-4})$ and
$1.40\times10^{-4}$.  Each arm completed all 20,000 GL2 time steps, and both
coupled implicit stage equations converged at every step.  The largest
delivered-$P_{128}$ Hamiltonian drift among the $R=10$
arms was $2.99\times10^{-8}$; this value is reported only as a numerical
diagnostic and is not a production acceptance condition.  We therefore use
the round margin $R\geq10$ and the exact two-times internal band
$K_{\mathrm{evol}}=256$ in Tanaka revision 3.

\paragraph{Equation-(33) Benjamin--Feir states.}
This family is restricted to deep water.  Let $n_c\in\mathbb N$ be the
carrier-mode index, let $\Delta n\in\mathbb N$ be the symmetric sideband
offset with $\Delta n<n_c$, let $\varepsilon_c>0$ be the steepness of the
first elevation harmonic, let $\rho>0$ be the sideband-to-carrier amplitude
ratio, and let $x_0\in[0,L)$ be a global spatial translation.  Thus
$\varepsilon_c$ is not defined from half the crest-to-trough wave height.
Define
\begin{equation}
  k_c=\frac{2\pi n_c}{L},\qquad
  k_-=\frac{2\pi(n_c-\Delta n)}{L},\qquad
  k_+=\frac{2\pi(n_c+\Delta n)}{L},\qquad
  a=\frac{\varepsilon_c}{k_c}.
  \label{eq:paper-bf-wavenumbers}
\end{equation}
The sideband ansatz follows equation~(33) of Xu and Guyenne, who in turn
follow Dommermuth and Yue \cite{dommermuth1987,xu2009}.  Their published test
uses, in the present notation,
\begin{equation}
  (n_c,\Delta n,\varepsilon_c,\rho,\theta)
  =
  (9,2,0.13,0.10,-\pi/4),
  \qquad L=2\pi,\qquad g=1.
  \label{eq:paper-bf-published-tuple}
\end{equation}
It also uses a numerically computed exact steady Stokes carrier.  The
construction below is our randomized generalization over 66 integer
mode-pair categories, not a reproduction of that single parameter tuple.

Let $y=x-x_0$, and let $(\eta_c(y),\xi_c(y))$ denote the project's analytic
fifth-order deep-water Stokes formula, with carrier phase zero in this
translated coordinate and first elevation harmonic of amplitude $a$.  This
analytic fifth-order carrier is an approximation to the published exact
steady carrier; it is not the Xu--Guyenne carrier itself.  Both sidebands use
the fixed relative phase $\theta=-\pi/4$ from the cited calculation.

With $\eta_0$ evaluated before $\xi_0$, the implemented sideband construction
is
\begin{align}
  \eta_0(x)
  &=
  \eta_c(y)+\rho a\left[
    \cos(k_-y+\theta)+\cos(k_+y+\theta)
  \right],
  \notag\\
  \xi_0(x)
  &=
  \xi_c(y)+\rho a\left[
    \sqrt{\frac{g}{k_-}}\,e^{k_-\eta_0(x)}\sin(k_-y+\theta)
    +
    \sqrt{\frac{g}{k_+}}\,e^{k_+\eta_0(x)}\sin(k_+y+\theta)
  \right].
  \label{eq:paper-bf-state}
\end{align}
The spatial mean of $\xi_0$ is removed after
\cref{eq:paper-bf-state}.  Both sidebands use the same $\rho$ and $\theta$;
there are no empirical cross terms.  Translating the complete state changes
neither the relative phases nor the resonant quartet phase, which remains
$\theta+\theta-2(0)=-\pi/2$.  The numerical depth is
$h=5L/(2\pi)$, so the first nonzero wavenumber $k_1=2\pi/L$ satisfies
$k_1h=5$.

For a proposed pair, define the normalized sideband location and a
focused-envelope steepness proxy by
\begin{equation}
  \beta
  =\frac{\Delta n/n_c}{2\sqrt{2}\,\varepsilon_c},
  \qquad
  \varepsilon_f
  =\varepsilon_c\left(1+2\sqrt{1-\beta^2}\right).
  \label{eq:paper-bf-band}
\end{equation}
Benjamin and Feir give the historical deep-water instability
condition~\cite{benjaminfeir1967}.  In the present normalization,
equation~(4.6) of Andrade and Stuhlmeier gives $0<\beta<1$, while their
equations~(5.7)--(5.11) give the Akhmediev-breather correspondence and the
peak-envelope amplification factor
$1+2\sqrt{1-\beta^2}$~\cite{andrade2023}.  We use that factor only to define
the scalar population-screening proxy above.  Yang and Liu use
$\varepsilon_c=0.10$, five-percent sidebands, and the maximum-growth
frequency coordinate $\Delta\omega/(\omega_c\varepsilon_c)=1$
\cite[Sec.~4.3.1]{yangliu2024}.  Under the leading deep-water narrowband
relation $\Delta k/k_c\simeq2\Delta\omega/\omega_c$, this corresponds to
$\beta\simeq1/\sqrt2$ in our symmetric-wavenumber normalization.  The
ten-percent endpoint below is the sideband ratio in the Xu--Guyenne test tuple in
\cref{eq:paper-bf-published-tuple}.  Revision 4 uses the declared support
\begin{equation}
  \begin{gathered}
    n_c\in\{4,\ldots,20\},\qquad
    \Delta n\in\{1,\ldots,n_c-1\},\qquad
    \varepsilon_c\in[0.05,0.13],\\
    0<\beta<1,qquad
    \varepsilon_f\leq F_0,qquad
    F_0=\frac{1+\sqrt2}{10},\\
    \rho\in[0.05,0.10],\qquad x_0\in[0,L).
  \end{gathered}
  \label{eq:paper-bf-support}
\end{equation}
The scalar $\varepsilon_f$ is a population-screening proxy, not a maximum
physical slope, a wave-breaking criterion, or a regularity theorem.

Let $\mathcal P$ be the 66 integer pairs $(n_c,\Delta n)$ for which this
support is nonempty.  Conditional on a pair in $\mathcal P$, set
\begin{equation}
  \ell=\frac{\Delta n}{2\sqrt2\,n_c},\qquad
  \varepsilon_{\min}=\max\{0.05,\ell\},\qquad
  \varepsilon_{\max}
  =\min\left\{0.13,
    \frac{-F_0+2\sqrt{F_0^2+3\ell^2}}{3}\right\}.
  \label{eq:paper-bf-conditional-steepness}
\end{equation}
The implemented sampler draws
$\varepsilon_c\sim\operatorname{Unif}
(\varepsilon_{\min},\varepsilon_{\max}]$,
$\rho\sim\operatorname{Unif}[0.05,0.10)$, and
$x_0\sim\operatorname{Unif}[0,L)$.  The upper endpoint in
\cref{eq:paper-bf-conditional-steepness} is obtained by solving
$\varepsilon_f=F_0$ at fixed $(n_c,\Delta n)$.  The relative phase
$\theta=-\pi/4$ is fixed rather than sampled.  The closed intervals in
\cref{eq:paper-bf-support} denote the closure of the physical support; the
half-open endpoint conventions above specify finite-precision sampling and do
not change the corresponding continuous uniform laws.
The quota rule treats each of the 66 pairs in
$\mathcal P$ as a parameter category.  If a family--split accepted-case quota is $C$, each
pair receives either $\lfloor C/66\rfloor$ or $\lceil C/66\rceil$ accepted
cases, in the fixed enumerated category order.  A rejected attempt is replaced in
the same pair category.  The pair-conditioned sampler and the general category
allocator are both implemented.

At the final split quotas, the 16384 training cases assign 248 or 249 accepted
cases to every pair, while the 1024 validation cases and 1024 test cases each
assign 15 or 16.  Additive training chunks contain the differences between
these cumulative balanced allocations.  Before bulk generation, the direct
revision-4 source gate accepted the first proposal in each of the 66
categories.  The completed production dataset subsequently retained 18432
cases from 18455 attempts: 16384 training cases, 1024 validation cases, and
1024 test cases.  The 23 rejected attempts retained no rows.  A source-bound
completion audit independently checked the ordered 66-category taxonomy and
quotas, exact proposal replay, transaction and array hashes, the 200
endpoint-pinned, approximately uniform selections from each saved grid,
stored depths and water columns, and the declared numerical conditions.  Its
artifact is
\path{outputs/paper_dataset_bf_revision4_jonswap_revision3_literature_aligned_v1/benjamin_feir_completion_audit.json},
with SHA-256
\par\noindent
\texttt{c155b35276d0cfef6844f6b0db3e91da}\allowbreak
\texttt{3ce15579ad8b0f0ca482f44cc912753c}.
Two shared Python modules changed after this calculation and their exact run
versions were absent from the Git history.  Those two files are now preserved,
without being imported by current code, in
\path{reproducibility/source_snapshots/}, in the subdirectory
\path{benjamin_feir_revision4_e16773f/};
the accompanying \texttt{SHA256SUMS} file equals the hashes recorded by every
one of the six chunk source maps.  The completion audit also reconstructs the
attempted, accepted, and rejected count in every category from the immutable
transactions and physically rehashes the other 17 recorded generation sources
against the current repository.  It remains the binding record for the dataset
arrays and transactions.
These counts measure the implemented population and numerical path; they do
not turn the proxy bound into a physical wave-breaking criterion.

The former common adapter persisted the pair, continuous parameters,
translation, and fixed relative phase before constructing the equation-(33)
state.  In the
reduced accepted-quota
calculation, the Benjamin--Feir attempt was accepted, committed three rows,
and was finite when read by the training loader.  This historical transaction
used the former paired validation path and exercised proposal, construction,
post-acceptance time selection, whole-case commit, and split-aware view
construction.  It checks order and storage logic only; the full-horizon
panel below is the separate method-level numerical test.

For each carrier, define the linear carrier period and the requested horizon
by
\begin{equation}
  T_c=\frac{2\pi}{\sqrt{gk_c}},\qquad
  H_{\mathrm{BF}}=100T_c,
  \qquad
  T_{\mathrm{BF}}
  =0.08\left\lfloor\frac{H_{\mathrm{BF}}}{0.08}\right\rfloor.
  \label{eq:paper-bf-horizon}
\end{equation}
The intended horizon is $H_{\mathrm{BF}}=100T_c$, and the calculation ends
at $T_{\mathrm{BF}}$, the last saved-grid time not exceeding that horizon.
This follows the carrier-period time coordinate used by Xu and Guyenne while
placing every endpoint on the common saved-time grid.  Across the 66 equally
weighted pair categories, $T_{\mathrm{BF}}$ ranges from $140.48$ to $314.08$
and has mean $174.50$; equivalently,
$T_{\mathrm{BF}}/T_c\in[99.9541,99.9973]$.

\paragraph{Resolved-band JONSWAP/TMA random seas.}
JONSWAP denotes the Joint North Sea Wave Project frequency spectrum, and TMA
denotes its standard Texel--Marsen--Arsloe finite-depth
modification.  Let $h>0$ be depth, let $k_p>0$ be the reference peak
wavenumber, let $\gamma\geq1$ be the peak-enhancement factor, and let $g>0$
be gravitational acceleration.  For wavenumber $k>0$, define
\begin{equation}
  \omega(k,h)=\sqrt{gk\tanh(kh)},\qquad
  c_g(k,h)=\frac{\partial\omega}{\partial k}
  =
  \frac{g[\tanh(kh)+kh\,\operatorname{sech}^2(kh)]}
       {2\omega(k,h)}.
  \label{eq:paper-random-dispersion}
\end{equation}
Set $\omega_p=\omega(k_p,h)$ and define the peak period
$T_p=2\pi/\omega_p$.  The intended terminal time, saved-time spacing, final
saved index, and realized terminal time are, respectively,
\begin{equation}
  H=16T_p=16\frac{2\pi}{\omega_p},
  \qquad
  \delta t_s=0.08,
  \qquad
  J=\left\lfloor\frac{H}{\delta t_s}\right\rfloor,
  \qquad
  T_s=J\delta t_s
      =0.08\left\lfloor\frac{H}{0.08}\right\rfloor.
  \label{eq:paper-random-horizon}
\end{equation}
Thus \(T_s\) is the last point on the saved-time grid that does not exceed
the intended horizon \(H\).  For angular frequency $\omega>0$, define
\begin{equation}
  \sigma(\omega)=
  \begin{cases}
    0.07,&\omega\leq\omega_p,\\
    0.09,&\omega>\omega_p,
  \end{cases}
  \label{eq:paper-jonswap-width}
\end{equation}
and the unnormalized JONSWAP density \cite{hasselmann1973}:
\begin{equation}
  \widetilde S_J(\omega)
  =
  g^2\omega^{-5}
  \exp\!\left[-\frac54\left(\frac{\omega_p}{\omega}\right)^4\right]
  \gamma^{\exp[-(\omega-\omega_p)^2/
    (2\sigma(\omega)^2\omega_p^2)]}.
  \label{eq:paper-jonswap-density}
\end{equation}
For the dimensionless number $z=kh$, define the TMA multiplier
\cite{bouws1985}:
\begin{equation}
  \Phi_{\mathrm{TMA}}(z)
  =
  \frac{\tanh^2z}{1+2z/\sinh(2z)}.
  \label{eq:paper-tma-factor}
\end{equation}
The cited works supply these spectral forms.  Published numerical JONSWAP
realizations have used the interval
$0.5f_p\leq f\leq2.5f_p$ \cite{draycott2026}, and an earlier random-wave
calculation used the upper endpoint
$\omega_{\max}=2.5\omega_p$ \cite{grice2015}.  We adopt this finite component
interval as a reproducible numerical definition, not as a wave-breaking
criterion.  The resolved Fourier cells, direction split, phase sampling, and
parameter strata below remain choices of the present dataset.
The finite-depth frequency density evaluated at wavenumber $k$ is
$\widetilde S_J(\omega(k,h))\Phi_{\mathrm{TMA}}(kh)$.  The factor
$c_g(k,h)=\partial\omega/\partial k$ in
\cref{eq:paper-random-dispersion} converts this frequency density to a
wavenumber density.

Let $\Delta k=2\pi/L$, let $k_j=j\Delta k$ for positive integers $j$, and
let $\mathcal J=\{j\in\mathbb N:k_j\leq K\}$ with $K=128$.  Revision 4
retains the fixed relative-frequency interval
\begin{equation}
  I_p(k)
  =
  \begin{cases}
    1,&\displaystyle \frac12\leq
       \frac{\omega(k,h)}{\omega_p}\leq\frac52,\\
    0,&\text{otherwise}.
  \end{cases}
  \label{eq:paper-random-relative-frequency-window}
\end{equation}
This is a sharp restriction in frequency relative to the prescribed peak; it
is not a cosine taper in absolute wavenumber.  The support requires
\begin{equation}
  \omega(K,h)\geq\frac52\,\omega_p,
  \qquad K=128,
  \label{eq:paper-random-relative-frequency-fit}
\end{equation}
so the complete upper end of the declared interval lies in the delivered
band for every proposed case.
For $j\in\mathcal J$, define the Fourier-cell endpoints
$a_j=\max\{0,k_j-\Delta k/2\}$ and
$b_j=\min\{K,k_j+\Delta k/2\}$.  Define the cell energy $E_j$ and normalized
energy fraction $\varpi_j$ by
\begin{align}
  E_j
  &=
  \int_{a_j}^{b_j}
  \widetilde S_J(\omega(k,h))
  \Phi_{\mathrm{TMA}}(kh)c_g(k,h)I_p(k)\,\dd k,
  \notag\\
  \varpi_j
  &=
  \frac{E_j}{\sum_{\ell\in\mathcal J}E_\ell}.
  \label{eq:paper-random-cell-energy}
\end{align}
Thus $\varpi_j\geq0$ and $\sum_{j\in\mathcal J}\varpi_j=1$.
The implementation evaluates each integral $E_j$ by 16-point
Gauss--Legendre quadrature.

Let $H_s>0$ be the ensemble significant wave height, define
$m_0=(H_s/4)^2$, and let $r_d\in[0,1]$ be the fraction of the linear wave
energy assigned to right-moving modes.  For every $j\in\mathcal J$, define
\begin{equation}
  A_{j,+}=\sqrt{2r_d m_0\varpi_j},\qquad
  A_{j,-}=\sqrt{2(1-r_d)m_0\varpi_j}.
  \label{eq:paper-random-amplitudes}
\end{equation}
The phase arrays $(\phi_{j,+})_{j\in\mathcal J}$ and
$(\phi_{j,-})_{j\in\mathcal J}$ are independent, and every entry is uniform
on $[0,2\pi)$.  Let $\G_0(h)$ be the flat-surface DNO, whose symbol at positive
wavenumber $k$ is $\G_0(k;h)=k\tanh(kh)$.  The initial state is
\begin{align}
  \eta_0(x)
  &=
  \sum_{j\in\mathcal J}\left[
    A_{j,+}\cos(k_j x+\phi_{j,+})
    +A_{j,-}\cos(k_j x+\phi_{j,-})
  \right],
  \notag\\
  \xi_0(x)
  &=
  \sum_{j\in\mathcal J}
  \frac{\omega(k_j,h)}{\G_0(k_j;h)}
  \left[
    A_{j,+}\sin(k_j x+\phi_{j,+})
    -A_{j,-}\sin(k_j x+\phi_{j,-})
  \right].
  \label{eq:paper-random-state}
\end{align}
The minus sign in the left-moving part of $\xi_0$ is the linear
traveling-wave relation.  The construction fixes
$\frac12\int_0^L(\xi_0\G_0(h)\xi_0+g\eta_0^2)\,\dd x=gLm_0$ independently
of phase; it does not rescale an individual realization to force its sampled
surface variance to equal $m_0$.

For every stratum, define the input peak steepness
\begin{equation}
  \varepsilon_p=\frac{k_pH_s}{2}.
  \label{eq:paper-random-peak-steepness}
\end{equation}
The declared support has three depth strata.  For the shallow stratum, let
$n_p\in\{16,18,20,22,24\}$, set $k_p=2\pi n_p/L$, and define
$\chi_p=k_p h$ and $\delta_s=H_s/(2h)$.  Its support is
\begin{equation}
  \chi_p\in[0.2,1.5],\qquad
  h=\chi_p/k_p,\qquad
  \delta_s\in[0.03,0.16],\qquad
  \varepsilon_p=\chi_p\delta_s\leq0.08.
  \label{eq:paper-random-shallow-support}
\end{equation}
The finite-depth and deep supports are, respectively,
\begin{align}
  k_p&\in[2,12],&
  h&\in[0.1,1.5],&
  H_s&\in[0.005,0.03],&
  \varepsilon_p&\leq0.08,
  \label{eq:paper-random-finite-support}\\
  k_p&\in[2,12],&
  h&\in[5,25],&
  H_s&\in[0.005,0.03],&
  \varepsilon_p&\leq0.08.
  \label{eq:paper-random-deep-support}
\end{align}
Every depth stratum also satisfies the resolved-frequency condition
\cref{eq:paper-random-relative-frequency-fit}.
Every stratum uses
$\gamma\in\{1,3.3,5\}$ and $r_d\in\{0,\tfrac12,1\}$.  The
parameter categories are the 27 combinations of depth stratum, $\gamma$, and
$r_d$, with requested accepted-case counts divided across those categories by
quotient and remainder.

The proposal law inside a shallow category is the following.  Draw
$n_p$ uniformly from the five values
\[
  n_p\in\{16,18,20,22,24\}.
\]
Draw $\chi_p$ uniformly on
$[0.2,1.5]$ and $\delta_s$ uniformly on $[0.03,0.16]$, repeating the pair
until both $\chi_p\delta_s\leq0.08$ and
\cref{eq:paper-random-relative-frequency-fit} hold.  This is the uniform
distribution, with respect to area, on the part of the displayed rectangle
that satisfies both conditions.  Then set $k_p=2\pi n_p/L$,
$h=\chi_p/k_p$, and $H_s=2h\delta_s$.  Inside a
finite-depth or deep category, draw $k_p$, $h$, and $H_s$ independently and
uniformly over the intervals in
\cref{eq:paper-random-finite-support} and
\cref{eq:paper-random-deep-support}, repeating the triple until
$k_pH_s/2\leq0.08$ and
\cref{eq:paper-random-relative-frequency-fit} both hold.  Equivalently, the
resulting proposed triple is uniform on the part of the corresponding
rectangular box that satisfies both displayed conditions.  In every
category, all entries of both phase arrays are independent and uniform on
$[0,2\pi)$.

The bound $\varepsilon_p\leq0.08$ is a restriction on the population being
studied, not a theorem that waves break immediately above it.  Ducrozet et
al. describe $k_pH_s/2=0.08$ as typical of a moderate JONSWAP sea state in a
HOS calculation~\cite{ducrozet2021rogue}.  A separate applicability study
shows that increasing the resolved HOS bandwidth can expose short-wave
content that a coarser calculation suppresses, thereby reducing the range in
which an apparently robust computation is accurate~\cite{ducrozet2017applicability}.
Since the present potential-flow solver has no physical wave-breaking model,
we use the moderate value as a conservative input restriction.  It does not
replace the full-trajectory numerical checks below: random phases and
nonlinear focusing can still make a proposed trajectory fail.

The spectrum, phase sampler, support predicate, state
\cref{eq:paper-random-state}, and 27-category proposal law are implemented.
Both realized phase arrays are held during construction but are not persisted;
they can be regenerated from the split, family, and attempt seed.  The adapter
then checks the
realized graph condition
\begin{equation}
  \min_\ell\{h+\eta_0(x_\ell)\}>0.
  \label{eq:paper-random-initial-graph-condition}
\end{equation}
This condition cannot be decided from $(h,H_s,k_p,\gamma,r_d)$ alone because
the minimum depends on the sampled phases.  If one realization fails, that
case contributes zero rows and is replaced in the same parameter category;
other cases in the same proposed numerical batch continue.  No deterministic
sum-of-amplitudes bound is used to reject otherwise valid phase realizations.

For JONSWAP/TMA, the event \(A_f\) in
\cref{eq:paper-accepted-law} specifically includes the phase-dependent initial
construction and graph condition above, the nonlinear adjustment, and the
autonomous trajectory decision.  This is not a separate phase-valued support
gate: it is the common complete-case conditioning specialized to this family.
The JONSWAP/TMA completion audit reports its effect separately; the analogous
common-law statement does not claim measured bias in another family.  No
retrospective parameter cutoff or bulk rejection-rate gate is applied.

The current adapter does not persist proposal values or construction metrics.
It connects a valid construction to the dedicated adjusted trajectory
executor, 16-time selection, accepted-row storage, and split-aware view.  An
earlier reduced accepted-quota calculation
used the former paired validation path: its random-sea attempt was accepted,
committed three finite rows, and was finite when read by the training loader.  Its
intended horizon was \(H=30.0341\), and its realized terminal time was
\(T_s=30.0\), in agreement with \cref{eq:paper-random-horizon}.  This is
historical software-path evidence, not a test of the current exact-contract
adjustment and autonomous path, the population law, or a population acceptance
rate.
The historical pre-adjustment 27-case pilot contains nine deterministic
support-spanning cases
per depth stratum and uses each $(\gamma,r_d)$ pair once per stratum; it was
not drawn from the population law and remains only a historical validation
panel.  It does not test the later adjustment and autonomous numerical
contract, nor does it test the revision-4 relative-frequency population.

Revision-4 JONSWAP/TMA generation does not insert the linear random sea
directly into the autonomous nonlinear equations.  Write the nonlinear
evolution equation as $z_t=Lz+R(z)$, where $z=(\eta,\xi)$, $Lz$ is its
linear gravity-wave part, and $R(z)$ is the nonlinear remainder.  During a
Dommermuth adjustment \cite{dommermuth2000}, the program instead solves
\begin{equation}
  z_t=Lz+A(t)R(z),
  \qquad
  A(t)=1-\exp\!\left[-\left(\frac{t}{10T_p}\right)^4\right].
  \label{eq:paper-jonswap-adjustment}
\end{equation}
Dommermuth supplies the ramped-nonlinearity construction; the exponent four,
the scale $10T_p$, and the $20T_p$ handoff horizon are fixed project
settings validated below rather than unique prescriptions of that method.
For each case, this nonautonomous calculation ends at
\begin{equation}
  T_{\mathrm{adj}}
  =0.08\left\lfloor\frac{20T_p}{0.08}\right\rfloor,
  \label{eq:paper-jonswap-adjustment-time}
\end{equation}
the last saved-grid time not exceeding $20T_p$.  Its complete internal-band
endpoint is passed to the autonomous solver without projection to
$P_{128}$ or componentwise rescaling.  The autonomous clock is then restarted
at $t=0$, $A$ is fixed to one, and the production interval has length
$T_s$ from \cref{eq:paper-random-horizon}.  Since
\cref{eq:paper-jonswap-adjustment} depends explicitly on time, Hamiltonian
conservation is not imposed during adjustment.  Finiteness, positive water
depth, and convergence of every implicit stage are still required.

\paragraph{Production rollout and whole-case acceptance.}
Tanaka cases are integrated over \(0\leq t\leq200\).  A Benjamin--Feir case
is integrated over \(0\leq t\leq T_{\mathrm{BF}}\), with
\(T_{\mathrm{BF}}\) defined in \cref{eq:paper-bf-horizon}; a random-sea
autonomous case is integrated over \(0\leq t\leq T_s\), with \(T_s\) defined
in \cref{eq:paper-random-horizon}.  Tanaka revision 3 uses
the order-six DNO on 1024 points in its filtered internal band
\(K_{\mathrm{evol}}=256\), as defined above.  Benjamin--Feir revision 4 uses
the order-four DNO on 1024 points in the sharp band \(|k|\le256\).
JONSWAP/TMA revision 4 uses the order-four DNO on 2048 points in the sharp
band \(|k|\le704\).  Both sharp-band methods use padding factor eight and
sharply project both state fields to \(P_{256}\) or \(P_{704}\), respectively,
after every completed GL2 step.  This repeated projection enforces the stated
Galerkin space but is not a Hou--Li filter: it applies no smooth damping to
retained modes.  Their constructed \(P_{128}\) states are embedded in the
corresponding internal grid and band; for JONSWAP/TMA the full-band result of
\cref{eq:paper-jonswap-adjustment} replaces that initial state before
autonomous evolution.  In every family the delivered state is a 1024-point
\((P_{128}\eta,P_{128}\xi)\) pair, and its label remains the order-six
quantity \(q_{\mathrm{ref}}\) in \cref{eq:forward-dataset-target}.  Thus the
order-four operator determines the Benjamin--Feir and random-sea
dynamics, whereas the order-six operator defines the learning label.

The production integrator uses GL2 step \(0.01\) and saves candidate states
every \(0.08\), without temporal interpolation.  JONSWAP/TMA permits at most
five fixed-point updates per implicit stage; the other trajectory contracts
retain their declared cap of four.  An attempted case must first
satisfy the parameter and construction support declared for its family.
Conditional on that support, a case is retained exactly when its declared
numerical path reaches the requested terminal time, every implicit stage is
marked converged, the stage and updated state are finite, every relative stage
residual is finite and at most \(10^{-8}\), all saved \(\eta\), \(\xi\), and
\(q_{\mathrm{ref}}\) values are finite, and
\[
  \min_{j,\ell}\{h+\eta(x_\ell,t_j)\}>0.
\]
Revision-4 Benjamin--Feir and JONSWAP/TMA cases have one further numerical
condition on their autonomous saved frames.  Let \(K_F=256\) for
Benjamin--Feir and \(K_F=704\) for JONSWAP/TMA, and define
\begin{equation}
  H_F(t_j)
  =\frac12\int_0^{2\pi}
    \left\{\xi\,G^{K_F}_4(\eta;h)\xi+g\eta^2\right\}\,\dd x,
  \qquad
  d_{H,F}
  =\max_j\frac{|H_F(t_j)-H_F(0)|}
                   {\max\{|H_F(0)|,\varepsilon_{64}\}},
  \label{eq:paper-internal-hamiltonian-gate}
\end{equation}
where \(G^{K_F}_4\) is the order-four DNO evaluated from the full sharp
internal state and \(\varepsilon_{64}\) is the smallest positive normal
float64 number.  The executor requires \(d_{H,F}\leq10^{-3}\).  Before
projection to \(P_{128}\), it also evaluates at every saved autonomous frame
whether the internal state and \(G^{K_F}_4(\eta;h)\xi\) are finite and the
minimum value of \(h+\eta\).  These dense saved-frame quantities, including
the initial Hamiltonian and maximum drift, are used for the acceptance
decision but are not persisted.
The Hamiltonian reference is the adjusted endpoint for JONSWAP/TMA; no
Hamiltonian gate is applied during its nonautonomous adjustment.

These statements define one requirement: the fixed method must return a
complete admissible trajectory.  The individual checks determine that
decision, but a rejected attempt's failure reason is not persisted.  They are
not separately fitted data filters.  An
incomplete case contributes no rows, including no finite prefix, and its
replacement is drawn from the same predeclared parameter category.  Sign
changes, slope, spectral roughness, target magnitude, visual appearance, and
model error do not enter this decision.  Hamiltonian drift enters only through
the declared autonomous revision-4 Benjamin--Feir and JONSWAP/TMA condition
\cref{eq:paper-internal-hamiltonian-gate}.

\paragraph{Method-level time-step validation.}
The production time step was fixed before bulk generation by a separate
comparison of two trajectories.  For steps $\Delta t_c$ and
$\Delta t_f=\Delta t_c/2$, let $y_{c,j}^{(\alpha)}$ and
$y_{f,j}^{(\alpha)}$ denote, at common saved time $t_j$, the coarse and fine
versions of the three fields
\begin{equation}
  y^{(1)}=\frac{P_K\eta}{h},
  \qquad
  y^{(2)}=\frac{P_K\xi}{h\sqrt{gh}},
  \qquad
  y^{(3)}=\frac{q_{\mathrm{ref}}}{\sqrt{gh}},
  \label{eq:paper-refinement-fields}
\end{equation}
where $g>0$ is gravitational acceleration.  With
$\lVert v\rVert_{2,N}^2=N^{-1}\sum_{\ell=0}^{N-1}|v_\ell|^2$, the refinement
defect is
\begin{equation}
  E
  =
  \max_{\alpha=1,2,3}
  \frac{\max_j
    \lVert y_{c,j}^{(\alpha)}-y_{f,j}^{(\alpha)}\rVert_{2,N}}
       {\max_j\lVert y_{f,j}^{(\alpha)}\rVert_{2,N}+10^{-12}}.
  \label{eq:paper-refinement-defect}
\end{equation}
The validation panel used $\Delta t_c=0.01$ and
$\Delta t_f=0.005$ and compared \(E\) with \(10^{-3}\).  This comparison,
and the optional \(0.005/0.0025\) diagnostic implemented by the audit utility,
are not run for individual production cases and do not choose which
trajectory is stored.

The two Stokes propagation cases in the predeclared full-horizon panel have
defects $E=5.42\times10^{-8}$ (finite depth) and
$8.24\times10^{-8}$ (deep water).  Five of the seven Tanaka and
Benjamin--Feir stress cases pass with $E\leq1.13\times10^{-5}$.  The remaining
steep seam-centered Tanaka case and equation-(33) stress case become
nonfinite at essentially the same physical times under both step sizes.
Thus halving the step changed no whole-trajectory decision.  Under the
production rule, those two incomplete \(0.01\) trajectories would contribute
no rows; their failure is not a large value of the refinement defect.  The
Tanaka members of this panel used the earlier $K_{\mathrm{evol}}=128$
contract, before revision 3 introduced the conditional support and buffered
evolution described above.  They are historical method-development evidence,
not a validation of the revision-3 Tanaka contract.

\paragraph{Historical pre-revision-3 failure-mechanism audit.}
Separate restart experiments diagnosed the two incomplete stress cases under
the earlier Tanaka and Benjamin--Feir contracts; they do not alter the dataset
rule.  The restart times are chosen before the
baseline tail exceeds a fixed guard, so we
distinguish the source-history measure from the post-restart arm measure.  Let
\(c_k^{\mathrm{src}}(t)\) be the \(k\)-th Fourier coefficient of the source
baseline DNO field, let \(c_k^{\mathrm r}(t)\) be that coefficient for one
restarted arm, and let \(t_r\) be its restart time.  In each arm, ``DNO field''
means \(P_K^\circ[\G_M^{N,p}(\eta(t);h)\xi(t)]\), evaluated with that arm's
values of \(N,M,p\), and \(K\).  For wavenumber bounds \(0\leq a\leq b\) and
any coefficient sequence \(c=(c_k)\), define
\[
  \mathcal E_{a:b}[c](t)=\sum_{a\leq |k|\leq b}|c_k(t)|^2,
  \qquad
  [z]_+=\max\{z,0\}.
\]
Here \(\mathcal E_{a:b}\) is the energy in wavenumbers
\(a\leq|k|\leq b\), and \([z]_+\) is the positive part of \(z\).  The two
growth measures are
\begin{align}
  G_{\mathrm B}^{\mathrm{src}}(t)
  &=
  \frac{
    [\mathcal E_{80:128}[c^{\mathrm{src}}](t)
     -\mathcal E_{80:128}[c^{\mathrm{src}}](0)]_+
  }{
    \mathcal E_{1:128}[c^{\mathrm{src}}](0)
  },
  \label{eq:paper-source-high-band-growth}\\
  G_{\mathrm B}^{\mathrm r}(t)
  &=
  \frac{
    [\mathcal E_{80:128}[c^{\mathrm r}](t)
     -\mathcal E_{80:128}[c^{\mathrm r}](t_r)]_+
  }{
    \mathcal E_{1:128}[c^{\mathrm r}](t_r)
  }.
  \label{eq:paper-high-band-growth}
\end{align}
The subscript ``B'' denotes the terminal band \(80\leq|k|\leq128\); the
superscripts ``src'' and ``r'' denote the source trajectory and restarted arm,
respectively.  At the selected restart, the source measures are
\(G_{\mathrm B}^{\mathrm{src}}(80)=4.65\times10^{-5}\) for Tanaka and
\(G_{\mathrm B}^{\mathrm{src}}(88)=6.26\times10^{-5}\) for
Benjamin--Feir, both below the predeclared \(10^{-4}\) guard.

In the one-factor Tanaka restart from \(t_r=80\), the
baseline \(N=1024,p=8\) arm and the \(N=1024,p=16\) arm both fail at
\(t=119.48\), whereas the \(N=2048,p=8\) arm remains finite with every GL2
stage solved through \(t=122\), but reaches
\(\max_t G_{\mathrm B}^{\mathrm r}(t)=0.562\).  The high-band threshold
times are nearly unchanged in all three arms.  A fourth arm changes only the
delivered cutoff
from \(K=128\) to \(K=192\).  It remains finite with every stage solved and
has \(\max_t G_{\mathrm B}^{\mathrm r}(t)=3.19\times10^{-5}\), below
\(10^{-4}\), in the former terminal band
\(80\leq |k|\leq128\), while permitting nonzero energy in
\(129\leq |k|\leq192\).  The largest energy in that added band is
\(2.80\times10^{-4}\) times the restart-time nonzero-mode energy, so the
improvement is not caused by a large amount of newly admitted energy.
The \(M=5\) and \(M=4\) arms also remain finite with every stage solved through
\(t=122\), but reach \(\max_t G_{\mathrm B}^{\mathrm r}(t)=0.463\) and
\(0.390\), respectively.
Lowering the series order therefore avoids the terminal solve failure in this
restart without removing the preceding high-band accumulation.

The corresponding Benjamin--Feir restart gives the opposite cutoff result.
The \(K=128,M=6\) baseline fails at \(t=108.78\), while changing only to
\(K=192\) makes the first failed stage earlier, at \(t=103.755\).  Changing
only to \(M=5\) remains finite with every stage solved through \(t=112\), but
still reaches \(\max_t G_{\mathrm B}^{\mathrm r}(t)=0.182\).  The \(M=4\)
arm also remains finite with every stage solved through \(t=112\), but reaches
\(\max_t G_{\mathrm B}^{\mathrm r}(t)=0.283\).  Over the finite saved prefix
of the \(K=192\) arm, its largest energy in \(129\leq|k|\leq192\), divided by
\(\mathcal E_{1:128}[c^{\mathrm r}](t_r)\), is \(0.0482\).  Thus extra product
padding does not remove either failure, a finer base grid prevents only the
Tanaka terminal breakdown, and a larger delivered band suppresses the Tanaka
accumulation but worsens the Benjamin--Feir failure.  Lower series order
avoids both terminal failures without suppressing their high-band growth:
both lower-order arms cross \(10^{-4}\), \(10^{-3}\), and \(10^{-2}\) at
\(t=92.32\), \(96.8\), and \(101.2\), the same times as the \(M=6\)
baseline.  The restart calculations inherit the baseline history before their
restart times and do not justify a change to the common dataset target.  Any
replacement numerical contract must instead pass a complete-horizon
comparison from the original initial condition.  The Tanaka revision-3
boundary panel summarized by \cref{eq:paper-tanaka-cross-resolution-error}
and \cref{eq:paper-tanaka-hidden-band-closure} is that separate comparison;
no restarted arm enters the revision-3 support decision.

In the historical pre-adjustment revision-2 panel, the shallow JONSWAP/TMA
case reaches the last saved-grid time not exceeding
16 peak periods with
$E=4.52\times10^{-6}$ and every implicit stage solved.  The finite-depth case
also reaches that last saved-grid time, with $E=5.63\times10^{-6}$, all 3864
coarse and 7728 fine stages solved, and no iteration-cap or nonfinite event.
The deep case does the same with $E=4.75\times10^{-6}$, all 3496 coarse and
6992 fine stages solved, and no cap or nonfinite event.  In total, 10 of the
12 fixed cases complete at both steps: both Stokes cases, three of four Tanaka
cases, two of three Benjamin--Feir cases, and all three random-sea cases.  No
\(0.0025\) diagnostic was invoked.  These are method-validation cases; they
do not constitute a generated dataset or estimate population acceptance
rates.

\paragraph{Historical JONSWAP/TMA spatial-cutoff validation.}
Before revision 4 changed the spectral population, the random-sea internal
cutoff was selected on a predeclared shallow revision-3 sentinel after its
complete nonlinear adjustment.  Both arms used 2048 internal grid
points, \(M=4\), padding factor eight, GL2 step \(0.01\), saved spacing
\(0.08\), residual tolerance \(10^{-8}\), and an iteration cap of five.  Over
the complete 16-peak-period autonomous history, the largest relative
whole-field discrepancies between independently evolved \(K=704\) and
\(K=768\) arms were
\begin{equation}
  (e_\eta,e_\xi,e_{G\xi})
  =(4.64267,1.93035,7.91931)\times10^{-4}.
  \label{eq:paper-jonswap-final-cutoff-errors}
\end{equation}
Here each field was transferred by the normalized Fourier construction in
\cref{eq:paper-normalized-fourier-coefficient}, and the order-six target was
recomputed on the common 1024-point grid.  The corresponding \(K=640\) versus
\(K=768\) calculation exceeded \(10^{-3}\).  Errors computed only from modes
\(96\le |k|\le128\) were about \(3\times10^{-3}\), so
\cref{eq:paper-jonswap-final-cutoff-errors} certifies only the three global
whole-field quantities, not a bandwise error bound.  In particular, this
calculation supports the inherited \(K=704\) evolution method but does not
validate the revision-4 relative-frequency population.

The then-current revision-3 source was replayed from the archived full
internal \(K=704\) endpoint.  It reproduced all 289 autonomous saved frames
with maximum relative \(L^2\) errors
\(3.99\times10^{-16}\), \(4.66\times10^{-16}\), and
\(1.17\times10^{-15}\) in \(\eta\), \(\xi\), and the recomputed target.  The
replay was accepted, solved every stage, had maximum residual
\(9.994198\times10^{-9}\), internal Hamiltonian drift
\(1.831187\times10^{-5}\), minimum water column \(0.0387699\), and CPU wall
time 197.355 seconds.  This verifies that the then-current source reproduces
the archived sentinel.  It remains historical method evidence.

The direct revision-4 population gate instead fixed 20 accepted cases in each
of the 27 categories before execution.  It retained all 540 requested cases
from 548 attempts.  Eight attempts failed the declared complete numerical
path and retained zero rows, giving rejection rate \(1.46\%\).  This is the
current-constructor gate used before revision-4 bulk generation.

\subsection{Historical v9 quality controls and the Tanaka generator revision}

The historical v9 acceptance tests are source-dependent because the dataset
was assembled in stages.  We therefore distinguish prevention during construction, rowwise
acceptance at assembly, whole-trajectory acceptance for the later nonlinear
additions, and an independent grid-refinement audit.  The sign-transition
screen used for an earlier Tanaka pilot dataset is documented separately in
\cref{app:historical-cleaning}; it was not applied to the v9 dataset used by the
reported model.

\paragraph{Prevention during construction.}
The first Tanaka pilot generator embedded a solitary-wave profile in the
periodic interval by interpolation with zero values outside the tabulated
profile.  If a randomly translated crest lay near an endpoint, this operation
cut off a non-negligible tail.  The resulting endpoint mismatch has slowly
decaying Fourier coefficients, which are amplified by the derivatives in a
high-order Craig--Sulem expansion.  The replacement Tanaka constructor removes
this mechanism before integration.  Each tabulated profile is extended by zero
outside its finite interval, and the generator sums every periodic image whose
support intersects $[0,L)$.  It evaluates each image by cubic Hermite
interpolation using the tabulated elevation and surface slope, evaluates the
sum on an $8N$ grid, and Fourier-restricts it to the $N$-point grid.  The
initial state and each rollout-derived target are then projected to
$|k|\leq128$.  Thus a profile that crosses $x=0$ is represented by periodic
wrapping rather than by truncation.

\paragraph{Historical row and case acceptance.}
The common assembly pass for the first ten archived sources retains a snapshot
only if
\begin{equation}
  \eta,\xi,q^\star,h,t\ \text{are finite},
  \qquad \norm{q^\star}_{\infty}\leq1,
  \label{eq:row-acceptance}
\end{equation}
where $q^\star$ denotes that source's stored target.
This pass removed \num{14751} of \num{2048000} Benjamin--Feir rows:
\num{14750} were nonfinite and one finite row violated the target-magnitude
bound.  Before assembly, the two random-sea archives had respectively 26 and
34 nonfinite rows removed.  No row from the two principal Tanaka archives was
removed by \cref{eq:row-acceptance}.

In historical v9 generation, the later shallow--steep and steep-Tanaka sources
use a stronger, case-level test.  At stored time $t_j$ and node $x_n$, let
$\eta_{j,n}$, $\xi_{j,n}$, and $q^\star_{j,n}$ denote the elevation,
potential, and source-specific stored target.  Define
\begin{equation}
  H_j=\frac{\Delta x}{2}\sum_{n=0}^{N-1}
      \left(\xi_{j,n}q^\star_{j,n}+g\eta_{j,n}^{2}\right),
  \qquad
  D_H=\max_j\frac{|H_j-H_0|}{|H_0|+10^{-30}}.
  \label{eq:data-hamiltonian-drift}
\end{equation}
A case is accepted only when all stored fields are finite,
$\min_{j,n}\eta_{j,n}>-0.9h$, and $D_H\leq\tau_H$.  We use
$\tau_H=10^{-5}$ for the three shallow--steep sources and
$\tau_H=10^{-3}$ for the steep-Tanaka addition.  The three shallow--steep
generation runs rejected 253, 230, and 209 cases, while the two steep-Tanaka
shards rejected 985 and 996 cases; generation continued until the target row
counts were reached.  These are acceptance tests on the stored temporal
subsample, rather than a claim that one identical filter was retroactively
applied to every source family.

\paragraph{Independent fixed-band resolution audit.}
Let $\vartheta$ denote the latent parameters from which an initial condition is
constructed.  We construct the same $\vartheta$ independently, in float64, on
$n_j=2^j N$ points for $j=0,1,2$; Fourier interpolation of the $N$-point state is
not an independent test.  Fix the delivered physical band $K$ and define
\begin{equation}
  q_j^K=P_K^\circ\!\left[
    \G_{6}^{n_j,8}(P_K\eta_{n_j};h)(P_K\xi_{n_j})
  \right].
  \label{eq:fixed-band-label}
\end{equation}
Let $Q_j^K$ be the vector of consistently normalized Fourier coefficients of
$q_j^K$ for $|k|\leq K$.  Our three-grid discrepancy is
\begin{equation}
  \Delta_K(\vartheta)
  =\max_{0\leq r<s\leq2}
  \frac{\norm{Q_r^K-Q_s^K}_{\ell^2}}
       {\norm{Q_s^K}_{\ell^2}+10^{-30}}.
  \label{eq:three-grid-discrepancy}
\end{equation}
The additive constant only defines the quotient for a vanishing target; it is
negligible for every case in the calibration panel.  We declare a relative
label accuracy budget $\tau_{\mathrm{ref}}=10^{-3}$ and require
$\Delta_K\leq\tau_{\mathrm{ref}}$ when accepting a generator configuration.

A matched Tanaka calibration used $N=1024$, $2N=2048$, $4N=4096$, and the
historical delivered band $K=341$.  It contained all 12 reconstructable
frame-zero sign rejections, four interior controls, and the same 12 rejected
crest specifications rebuilt with the corrected periodic construction.
\Cref{tab:three-grid-calibration} shows the resulting separation.  Decisions
were unchanged when
$\tau_{\mathrm{ref}}$ varied over
$3\times10^{-4}$, $10^{-3}$, and $3\times10^{-3}$.  For the matched periodic
counterfactuals, the median $N$--$2N$ label defect was
$7.03\times10^{-6}$ and the median $2N$--$4N$ defect was
$2.18\times10^{-6}$; the clipped construction gave medians $0.0827$ and
$0.147$, respectively.  This is a convergence test for the fixed order-six
target on a fixed physical band.  It does not test Craig--Sulem order
truncation, physical admissibility, or time-integration error, which are
addressed separately.

\begin{table}[t]
  \centering
  \caption{Three-grid fixed-band discrepancy on the matched Tanaka
  calibration.  The same decision is obtained at all three tested tolerances.}
  \label{tab:three-grid-calibration}
  \begin{tabular}{lrr}
    \toprule
    Construction & Cases & Range of $\Delta_K$\\
    \midrule
    Corrected periodic counterfactual & 12 &
      $5.86\times10^{-6}$--$1.10\times10^{-5}$\\
    Historical interior control & 4 &
      $1.46\times10^{-4}$--$1.64\times10^{-4}$\\
    Historical clipped profile & 12 &
      $2.33\times10^{-2}$--$7.08\times10^{-1}$\\
    \bottomrule
  \end{tabular}
\end{table}

For the existing v9 archive this is a generator-level resolution audit, not a
claim that all \num{7633989} snapshots were retrospectively filtered by
\cref{eq:three-grid-discrepancy}.  For the replacement dataset, this particular
three-grid statistic audits the Tanaka construction.  The separate cutoff
comparisons for other families have only the scope stated with each panel; in
particular, the historical revision-3 JONSWAP/TMA sentinel does not validate
the revision-4 input population.  None of these method-level comparisons is
repeated as a per-candidate state or label filter.  Individual rollout cases
are governed by the declared-family and whole-trajectory rules above.

\paragraph{Historical pilot cleaning.}
The April Tanaka pilot dataset was in fact cleaned with a dead-zoned sign-count
rule before its 500,000-row training subset was drawn.  We retain that fact for
reproducibility, but do not present the count as a general numerical validity
criterion.  Its exact definition, removal count, mechanism, and matched
periodization counterfactual are given in \cref{app:historical-cleaning}.

\subsection{Sampling, splitting, and normalization}

\paragraph{Historical v9 split.}
The historical v9 archive has no trajectory map used by the training loader.
If $R$ denotes its total number of stored rows, the loader applies a
fixed-seed random permutation to the integers $0,\ldots,R-1$.  The validation
and test counts are each $\lfloor0.1R\rfloor$, and all remaining rows are
training rows.  The training rows occur first in the permutation, followed
by validation and then test.  This is a row split: two saved times from the same
trajectory can occur in different subsets.  That fact must be retained when
interpreting the reported v9 holdout error.

\paragraph{Replacement-dataset temporal selection.}
Temporal selection is applied only after a complete reference trajectory has
passed the numerical acceptance calculation.  A Tanaka trajectory has
candidate times $t_j=0.08j$, $j=0,\ldots,2500$.  Define
\begin{equation}
  S_{\mathrm T,j}
  =
  \sum_{\ell=0}^{N-1}
  |\partial_x\eta(x_\ell,t_j)|^2,
  \qquad
  A_{\mathrm T,j}
  =
  \frac{|D S_{\mathrm T,j}|}{S_{\mathrm T,j}+10^{-12}},
  \label{eq:paper-tanaka-activity}
\end{equation}
where $D S_{\mathrm T,j}=(S_{\mathrm T,j+1}-S_{\mathrm T,j-1})/2$ at
interior indices and is the corresponding one-sided difference at the two
endpoints.  This activity is smoothed by a normalized discrete Gaussian with
standard deviation 50 candidate-time indices.  Values beyond either endpoint
are set equal to the nearest endpoint value during this convolution.  Denote
the smoothed sequence by $\widetilde A_{\mathrm T,j}$.  The sampling mass is
\begin{equation}
  p_{\mathrm T,j}
  =
  \frac{1-\alpha_{\mathrm T}}{2501}
  +
  \alpha_{\mathrm T}
  \frac{\widetilde A_{\mathrm T,j}}
       {\sum_{r=0}^{2500}\widetilde A_{\mathrm T,r}},
  \qquad \alpha_{\mathrm T}=0.5.
  \label{eq:paper-time-density}
\end{equation}
If the smoothed activity is identically zero, its normalized term in
\cref{eq:paper-time-density} is replaced by the uniform mass $1/2501$.
A deterministic inverse-CDF rule first computes indices at the
midpoint quantiles $(\nu+\tfrac12)/200$, where
$\nu=0,\ldots,199$.  The first and last indices are fixed to 0 and 2500.
Each remaining index is restricted so that enough later integer indices
remain available and is then advanced if necessary.  This gives exactly 200
strictly increasing indices even when the activity density is concentrated.

For a Benjamin--Feir trajectory, let
$t_j=0.08j$, $j=0,\ldots,J_{\mathrm{BF}}$, where
$J_{\mathrm{BF}}=T_{\mathrm{BF}}/0.08$.  Retain the 200 indices
\begin{equation}
  j_\ell
  =\left\lfloor\frac{\ell J_{\mathrm{BF}}}{199}+\frac12\right\rfloor,
  \qquad \ell=0,\ldots,199.
  \label{eq:paper-bf-retained-indices}
\end{equation}
Thus the realized interval $[0,T_{\mathrm{BF}}]$, which is approximately
100 carrier periods, is sampled approximately uniformly, including both
endpoints, without introducing an additional event score.

For a random-sea trajectory, let \(t_j=j\delta t_s\),
\(j=0,\ldots,J\), be the dense saved times, where \(\delta t_s=0.08\) and
\(J=\lfloor H/\delta t_s\rfloor\) are defined in
\cref{eq:paper-random-horizon}.  Here \(t=0\) is the adjusted full-band
endpoint and the beginning of the autonomous production clock, not the
linear random sea before adjustment.  For each \(\ell=0,\ldots,15\), retain the
index
\begin{equation}
  j_\ell=
  \left\lfloor\frac{\ell J}{15}+\frac12\right\rfloor.
  \label{eq:paper-random-retained-indices}
\end{equation}
Thus \(j_\ell\) is the nearest integer to \(\ell J/15\), and the retained time
is \(t_{j_\ell}\).  In particular, \(t_{j_{15}}=T_s\) is the last saved-grid
time not exceeding \(H=16T_p\).  A static Stokes case retains only \(t=0\).
These rules change which accepted states are stored; they do not change the
target in \cref{eq:forward-dataset-target}.

\paragraph{Replacement-dataset case split and loader sampling.}
An attempted case is assigned to training, validation, or test before its
initial condition is constructed.  The three split roots are 2026072210,
2026072204, and 2026072205, respectively.  A rejection remains attached to
its assigned split and parameter category, and no saved times from one case may
belong to different splits.  In each epoch, the training loader uses every
accepted training row and globally reshuffles the row indices.  Equal accepted-
case counts across the four families therefore do not imply equal row weight
when the families retain different numbers of times per case.  Validation uses
every accepted validation row in one fixed deterministic order.

A training case from each of the four families contributes
\[
  1+200+200+16=417
\]
stored rows in total.  Training is generated in four cumulative stages,
\[
 C\in\{2048,4096,8192,16384\}
\]
accepted cases per family, while validation and test remain fixed at 1024
cases per family.  The first stage therefore contains
\(417(2048+1024+1024)=1{,}708{,}032\) rows and is a learning-curve checkpoint,
not the completed dataset.  The final target contains
\(417(16384+1024+1024)=7{,}686{,}144\) rows, which is 52,155 rows (0.68\%)
larger than the 7,633,989-row historical v9 dataset.  The four additive
training-stage increments contain 2048, 2048, 4096, and 8192 cases per family,
respectively; a completed chunk is never enlarged in place.  An increment may
itself be divided into disjoint execution shards.  Stokes, Tanaka, and
Benjamin--Feir use one shard for each of these four increments.  To balance
the two longer-running JONSWAP/TMA worker lanes, its final 8192-case increment
uses the three disjoint intervals \([8192,12288)\), \([12288,14336)\), and
\([14336,16384)\).

The split-tagged attempt allocator, atomic completed-batch storage, multi-shard
manifest, row map, and loader are implemented and tested as family-independent
modules.  Each completed batch stores its family and split, its ordered
parameter-group identifiers, and any accepted row arrays.  A rejected slot is
identified by the absence of owned rows.  The loader uses the stored split and
globally shuffles every accepted training row each epoch.  The historical
exact-target Stokes calculation described above exercised the former
transaction.  Under
Stokes revision 2, one case from each of its four parameter categories passed
this exact calculation, with all four quality-decision masks equal to zero.
Its completed production splits contain 16384 training cases and 1024 cases
in each of validation and test; these artifacts are retained.
In the former reduced accepted-quota workflow, the Tanaka, Benjamin--Feir, and
random-sea adapters enforced sampling, durable proposal, then construction.
One case from each of these three families continued through the former paired
GL2 validation path and loader-facing view, and each produced three finite
loaded rows.  Those
family-revision-1 calculations are historical transaction tests, not tests of
the current family-revision executors, numerical validation of the paper
target, or evidence about any family population.

The exact Tanaka revision-3 pilot was instead run with the paper contract in
two complementary validation-tagged shards covering categories 0--5 and
6--10.  All 11 attempted cases were accepted, no case was rejected, and the
11 trajectories produced 2200 rows.  These rows are diagnostic pilot output:
their dedicated roots are excluded from the replacement-dataset manifest,
training, and normalization despite their validation split tags.  A later
historical Tanaka run completed 4096 cases and retained 200 rows from each,
but a constructor audit found 64 mislabeled initial conditions caused by the
small-amplitude floor described above.  A corrected initial condition changes
the nonlinear trajectory and its selected times, so the old rows cannot be
repaired algebraically or mixed with corrected rows.  The corrected
revision-3 release instead draws a fresh deterministic population with 16384
training cases and 1024 cases in each of validation and test.  It explicitly
does not reuse the historical random specifications.

All four family releases and all four cumulative dataset views are complete.
The final release retained \num{73728} cases from \num{74045} attempts,
recorded \num{317} zero-row rejections, and stored \num{7686144} rows.  The
family totals (attempted/accepted/rejected/rows) are Stokes revision~2,
\num{18432}/\num{18432}/0/\num{18432}; Tanaka revision~3,
\num{18432}/\num{18432}/0/\num{3686400}; Benjamin--Feir revision~4,
\num{18455}/\num{18432}/23/\num{3686400}; and JONSWAP/TMA revision~4,
\num{18726}/\num{18432}/294/\num{294912}.  No reported model has been trained
on this replacement release.

\begin{table}[t]
  \centering
  \caption{Attempted/accepted replacement-dataset cases by family and split.}
  \label{tab:replacement-family-split-counts}
  \begin{tabular}{lrrr}
    \toprule
    Family & Training & Validation & Test\\
    \midrule
    Stokes & 16384/16384 & 1024/1024 & 1024/1024\\
    Tanaka & 16384/16384 & 1024/1024 & 1024/1024\\
    Benjamin--Feir & 16404/16384 & 1024/1024 & 1027/1024\\
    JONSWAP/TMA & 16644/16384 & 1040/1024 & 1042/1024\\
    \bottomrule
  \end{tabular}
\end{table}

The declared category marginals also close exactly.  Each of the four Stokes
categories has \num{4608}/\num{4608} attempted/accepted cases.  Across the
eleven Tanaka categories, six contain \num{1675}, four contain \num{1676}, and
one contains \num{1678} attempted and accepted cases.  Across the 66
Benjamin--Feir categories, 32 contain 278 accepted cases, 18 contain 280, and
16 contain 281; only \texttt{n\_c\_04\_\_delta\_n\_01} has rejections
(304 attempted, 281 accepted), while every other category has
attempted=accepted.  The 27 JONSWAP/TMA categories contain 22 accepted quotas
of 683, three of 682, and two of 680; 16 categories account for all 294
rejections.  The immutable document handoff records every named category's
exact attempted, accepted, and rejected counts rather than reducing those
JONSWAP/TMA counts to a bulk threshold.

The exact SHA-256 identities of the cumulative combined summaries, indexed by
training cases per family, are
\begin{description}
  \item[2048:] \texttt{2265cb71270accb3}\allowbreak
    \texttt{f6d9103e90433b664}\allowbreak
    \texttt{ca813989dfab6340}\allowbreak
    \texttt{3d8612af22d28c3};
  \item[4096:] \texttt{7feb018ee4d94eab}\allowbreak
    \texttt{10278917c1452620}\allowbreak
    \texttt{26cccab6169eaaa2}\allowbreak
    \texttt{d8877f7c7b28c9da};
  \item[8192:] \texttt{43b1eb91198e724c}\allowbreak
    \texttt{6f32b871c2d0efd0}\allowbreak
    \texttt{d8355533789919df}\allowbreak
    \texttt{5897ad588e7e85f6};
  \item[16384:] \texttt{32f764cf865c9089}\allowbreak
    \texttt{2ee0329e48fe992c}\allowbreak
    \texttt{bc24df2707e05fd1}\allowbreak
    \texttt{af2ee9cd2b84307c}.
\end{description}
The final summary binds the following 26-source manifest and trajectory-map
SHA-256 identities:
\begin{center}
\small
\texttt{dbb76ed1cba4d0146f3c9643a73d168d3a9fa97bc27c37fe0fd7ad08554bd734}\\
\texttt{f26888ad0a2c738405fb1662b9410337e1babfe48f7913805b5243ae0f279cc8}.
\end{center}

\paragraph{Historical and replacement-dataset normalization.}
Let $\mathcal I_{\mathrm{v9}}$ be the set of all historical v9 row indices.
Let $\eta_i$ and $\xi_i$ denote the two input fields in row $i$, and let
$q_i^\star$ denote its stored target under that source's historical label
convention.  The reported scale-normalized v9 checkpoints use the
dataset-wide absolute maxima
\begin{equation}
  s_\eta^{\mathrm{v9}}
  =\max_{i\in\mathcal I_{\mathrm{v9}},\,x}|\eta_i(x)|,\qquad
  s_\xi^{\mathrm{v9}}
  =\max_{i\in\mathcal I_{\mathrm{v9}},\,x}|\xi_i(x)|,\qquad
  s_q^{\mathrm{v9}}
  =\max_{i\in\mathcal I_{\mathrm{v9}},\,x}|q_i^\star(x)|.
  \label{eq:paper-v9-scales}
\end{equation}
For the replacement dataset, let $\mathcal I_{\mathrm{train}}$ contain only rows
from accepted cases preassigned to training.  For
$u\in\{\eta,\xi,q_{\mathrm{ref}}\}$, let $u_i(x)$ denote field $u$ at
coordinate $x$ in row $i$, and define
\begin{equation}
  s_u
  =
  \max_{i\in\mathcal I_{\mathrm{train}},\,x}|u_i(x)|,
  \qquad
  u\in\{\eta,\xi,q_{\mathrm{ref}}\}.
  \label{eq:paper-training-scales}
\end{equation}
For the authenticated release these training-only scales are
\(s_\eta=0.15583430230617523\),
\(s_\xi=0.09394174814224243\), and
\(s_{q_{\mathrm{ref}}}=0.12198150902986526\), computed from
\num{6832128} training rows.  Their statistics fingerprint is
\begin{center}
\small
\texttt{4e8ec96947bc05cee3b375467499b0ae00a7814f9982614084efd859efe7c3b8}.
\end{center}
For one replacement-dataset row, the scale-normalized fields are
\begin{equation}
  \widetilde\eta=\eta/s_\eta,\qquad
  \widetilde\xi=\xi/s_\xi,\qquad
  \widetilde q_{\mathrm{ref}}=q_{\mathrm{ref}}/s_{q_{\mathrm{ref}}}.
  \label{eq:normalization}
\end{equation}
The historical formula has the same form after replacing
$(s_\eta,s_\xi,s_{q_{\mathrm{ref}}},q_{\mathrm{ref}})$ by
$(s_\eta^{\mathrm{v9}},s_\xi^{\mathrm{v9}},s_q^{\mathrm{v9}},q^\star)$.
The replacement-dataset depth coordinate is
$d_h=\log(h/h_\star)$, where $h_\star=L/(2\pi)$ is a fixed reference length.
Because $L=2\pi$ in the present calculations, $d_h$ has the same numerical
value as the historical input $\log h$.  The analytic $\G_0+\G_1$ paths
recover physical $\eta$ and $\xi$ before applying their Fourier symbols, and
predictions are returned to physical target units for rollout.  This
empirical normalization is distinct from the exact coordinate dilation
derived in \cref{app:rescaling}: it changes neither the period nor the depth.

Several audits distinguish missing data from missing structure.  Random and
source-stratified scans found no nonfinite base rows, and models that later
fail in rollout can still predict the stored target $q^\star$ accurately on
the corresponding exact truth states.  Fine-tuning on states harvested from an
older model rollout was actively harmful.  Their order-six labels were correct
for the harvested states,
but those off-manifold inputs already contained high-band contamination created
by the older model; the resulting physical high-band response was therefore an
undesirable training signal, not a corrupted label.  These
observations motivate the shape-consistency constraint in
\cref{sec:hamiltonian-loss}: more pointwise samples on the same trajectory
manifold do not determine the derivative transverse to that manifold.
